% !TEX program = xelatex
% Converted from the current numbered Markdown draft sections, 18 October 2026.
% ESB Part IV template v2: layout and title-page commands retained.
\documentclass[12pt]{article}
\usepackage{ESB_preamble}
\usepackage{placeins}
\usepackage{booktabs,array,calc,amssymb,pdflscape}
\usepackage{seqsplit,xurl}
\urlstyle{same}
\newcommand{\tightlist}{\setlength{\itemsep}{0pt}\setlength{\parskip}{0pt}}
\newcommand{\pandocbounded}[1]{#1}
\newcounter{none}
\setcounter{secnumdepth}{3}
\setcounter{tocdepth}{2}
\setlength{\emergencystretch}{3em}
\setlength{\headheight}{14pt}
\renewcommand{\topfraction}{0.95}
\renewcommand{\textfraction}{0.05}
\newcommand{\ESBassessment}{Part IV Research Project Report}
\newcommand{\ESBtitle}{Development of a geothermal separator CFD model: liquid transport, collection and wall-film behaviour}
\newcommand{\ESBauthor}{Shuhei Yokkaichi}
\newcommand{\ESBcoworker}{Andy Tran}
\newcommand{\ESBsupervisors}{Sadiq Zarrouk \\ Maedeh Amirpour}
\newcommand{\ESBdate}{18 October 2026}
\newcommand{\ESBstudentid}{302558861 / syok443}
\hypersetup{pdftitle={Development of a geothermal separator CFD model},pdfauthor={Shuhei Yokkaichi}}

% Compact report layout: body font, line spacing and page margins remain ESB values.
\setlength{\parskip}{4pt plus 1pt minus 0.5pt}
% Keep a positive gap between headings and the following text at every level.
\titlespacing*{\section}{0pt}{6pt plus 2pt minus 1pt}{6pt plus 1pt minus 1pt}
\titlespacing*{\subsection}{0pt}{4pt plus 1pt minus 1pt}{6pt plus 1pt minus 1pt}
\titlespacing*{\subsubsection}{0pt}{4pt plus 1pt minus 1pt}{5pt plus 1pt minus 1pt}
\setlength{\textfloatsep}{10pt plus 2pt minus 2pt}
\setlength{\intextsep}{9pt plus 2pt minus 2pt}
\setlength{\floatsep}{9pt plus 2pt minus 2pt}
\captionsetup{skip=5pt}
\renewcommand{\bottomfraction}{0.9}
\renewcommand{\floatpagefraction}{0.8}
\setcounter{topnumber}{5}
\setcounter{bottomnumber}{3}
\setcounter{totalnumber}{8}
% Keep equations beside their introducing and defining text.
\predisplaypenalty=10000
\postdisplaypenalty=10000
\interdisplaylinepenalty=10000
\clubpenalty=10000
\widowpenalty=10000
\raggedbottom
\makeatletter
\setlength{\@fptop}{0pt}
\setlength{\@fpsep}{10pt plus 2pt minus 2pt}
\makeatother

\begin{document}
\pagenumbering{roman}
\hypersetup{pageanchor=false}
\input{ESB_titlepage.tex}
\hypersetup{pageanchor=true}

\ESBfrontsection{Abstract}
Geothermal steam-water separators must remove liquid efficiently while preventing its re-entrainment into the steam stream. To investigate liquid transport within a separator and support design improvements, this study extended an existing computational fluid dynamics (CFD) model. The original model lacked a dedicated liquid discharge pathway and did not explicitly represent wall-film behaviour. To overcome these limitations, modifications to the inlet conditions, mesh resolution, and startup procedure were combined with the implementation of liquid removal and wall-film modelling, thereby enabling detailed examination of these processes.

Liquid-distribution results revealed flow features near the separator wall that were not evident from outlet flow rates alone. In the absence of a liquid collection pathway, liquid continued to accumulate within the separator, demonstrating that low liquid carryover through the steam outlet does not necessarily indicate effective separation. Introducing liquid removal and initially reducing the inlet flow rate reduced the reported scaled continuity residual by approximately a factor of 650 and decreased the retained liquid inventory by about 93\% relative to an earlier case without liquid removal. Although direct comparison is complicated by differences in scaling, inlet flow conditions, and mesh resolution, the lower residual levels and minimal variation in inventory during the reduced-flow period provided a stable basis for investigating wall roughness, film transport, and drainage behaviour.

With the surrounding flow field held constant, wall-film evolution was simulated for 500 ms. During this period, the liquid film spread over a large portion of the separator wall. Although drainage rates approached a quasi-steady state, the total liquid inventory continued to increase. Removal of film near the separator base reduced local accumulation; however, liquid continued to collect in upper wall regions. Simulations using finer meshes produced more distinct drainage pathways, lower liquid inventories, and reduced drainage rates, indicating that the influence of mesh resolution remains unresolved.

These findings informed two design concepts aimed at improving liquid collection: a vertical barrier to interrupt the recirculation of wall liquid and an inclined ramp to direct liquid downward toward the collection zone. While the predicted flow patterns indicated favourable local changes, enhanced liquid collection has not yet been demonstrated. The duration of the simulations was constrained by computational cost; however, the upgraded model now provides a platform for more detailed investigation of liquid transport and collection processes. Future work should determine whether wall-film accumulation eventually reaches equilibrium and whether the proposed design modifications improve drainage performance. Before the model can be used confidently to predict separator performance, a complete liquid mass balance and validation against experimental measurements are required.

\clearpage

\ESBfrontsection{Declaration}
\textbf{Student}

I hereby declare that:

\begin{enumerate}
\def\labelenumi{\arabic{enumi}.}
\tightlist
\item
  This report is the result of the final year project work carried out by my project partner (see cover page) and I under the guidance of our supervisor (see cover page) in the 2026 academic year at the Department Engineering Science and Biomedical Engineering, Faculty of Engineering, University of Auckland.
\item
  This report is not the outcome of work done previously.
\item
  This report is not the outcome of work done in collaboration, except that with a potential project sponsor (if any) as stated in the text.
\item
  This report is not the same as any report, thesis, conference article or journal paper, or any other publication or unpublished work in any format.
\end{enumerate}

\textbf{Continuing-project statement.} The work builds on an inherited separator model and the published work of Purnanto et al.~(2013). The present investigation developed liquid-collection and startup procedures, examined mesh and inlet assumptions, and studied wall roughness, film transport, release and drainage. The starting geometry and bulk modelling approach are identified as inherited work.

\textbf{Use of generative AI.} OpenAI Codex assisted with literature retrieval, coding and checking Fluent/PyFluent workflows, execution and recovery support, analysis, figure preparation and report drafting. Fluent performed the numerical calculations. I directed the scientific scope and supplied the physical interpretation.

\textbf{Word count:} 9,817

\textbf{Signature and date:} Shuhei Yokkaichi, 18th of October 2026

\clearpage
\tableofcontents\clearpage\listoffigures\clearpage\listoftables\clearpage

\ESBfrontsection{Glossary of terms and abbreviations}
\begingroup\small\setlength{\tabcolsep}{3pt}
{\def\LTcaptype{none} % do not increment counter
\begin{longtable}[]{@{}
  >{\raggedright\arraybackslash}p{(\linewidth - 2\tabcolsep) * \real{0.5000}}
  >{\raggedright\arraybackslash}p{(\linewidth - 2\tabcolsep) * \real{0.5000}}@{}}
\toprule\noalign{}
\begin{minipage}[b]{\linewidth}\raggedright
Term
\end{minipage} & \begin{minipage}[b]{\linewidth}\raggedright
Meaning in this report
\end{minipage} \\
\midrule\noalign{}
\endhead
\bottomrule\noalign{}
\endlastfoot
BOC & Bottom-outlet cyclone \\
CFD & Computational fluid dynamics \\
Carrier & The Eulerian vapour and bulk-liquid flow fields through which discrete droplets are tracked \\
DPM & Discrete Phase Model; the Lagrangian representation of droplets \\
EWF & Eulerian Wall Film; the wall-based representation of liquid-film transport \\
IAPWS & International Association for the Properties of Water and Steam \\
IF97 & IAPWS industrial formulation for equilibrium water and steam properties \\
PyFluent & Python interface used to configure and command Ansys Fluent \\
Agent harness & Recorded workflow linking research questions, agent instructions, solver execution and analysis \\
RNG & Renormalisation group; used here in the RNG k--ε turbulence model \\
SIMPLE & Semi-Implicit Method for Pressure-Linked Equations \\
Bulk liquid & Liquid represented by the Eulerian secondary phase; excludes EWF inventory and discrete droplets \\
Vapour / bulk liquid & The primary vapour and secondary bulk-liquid phases in the selected Fluent carrier models \\
Accretion & Transfer of represented bulk liquid into the wall film \\
Stripping & Release of liquid from a film under gas forcing \\
Splashing & Production of secondary droplets during droplet impact on a film \\
Edge separation & Release of film liquid at a wall edge \\
Direct film drain & Local mass and momentum sink acting on lower-wall film storage \\
Incomplete trajectory & A tracked droplet with no completed terminal fate under the stated tracking controls \\
Absorber / virtual collector & A numerical liquid-removal treatment inside the truncated model; the two terms refer to the same model role \\
Native iteration & Fluent's recorded solver-update coordinate; it is not physical carrier-flow time \\
Film time & Time accumulated from the accepted EWF steps \\
Source term & A rate added to a transport equation; a negative mass source removes mass \\
Implicit source linearisation & Use of source derivatives within an equation solve to represent source dependence on the solved variable \\
Numerical stability & Observed finite and controlled behaviour over the stated calculation window \\
Iterative convergence & Active equations and relevant outputs meeting their declared settling conditions \\
Numerical accuracy & Accuracy of the computed solution of the stated model, assessed through numerical error and resolution evidence \\
Stiff source & A source with strong sensitivity to the solved variable, which can make an equation harder to solve \\
Numerical qualification & Assessment of a calculation against declared numerical, accounting and stationarity checks \\
Physical validation & Assessment against relevant physical observations for a stated use \\
\end{longtable}
}
\endgroup

\bigskip

\ESBfrontsection{Symbols}
\begingroup\small\setlength{\tabcolsep}{3pt}
{\def\LTcaptype{none} % do not increment counter
\begin{longtable}[]{@{}
  >{\raggedright\arraybackslash}p{(\linewidth - 4\tabcolsep) * \real{0.20}}
  >{\raggedright\arraybackslash}p{(\linewidth - 4\tabcolsep) * \real{0.60}}
  >{\raggedright\arraybackslash}p{(\linewidth - 4\tabcolsep) * \real{0.20}}@{}}
\toprule\noalign{}
\begin{minipage}[b]{\linewidth}\raggedright
Symbol
\end{minipage} & \begin{minipage}[b]{\linewidth}\raggedright
Meaning
\end{minipage} & \begin{minipage}[b]{\linewidth}\raggedright
Unit
\end{minipage} \\
\midrule\noalign{}
\endhead
\bottomrule\noalign{}
\endlastfoot
\(\dot m\) & Mass-flow rate & kg/s \\
\(M_b, M_f\) & Bulk-liquid and film inventories & kg \\
\(M_c\) & Liquid mass in the collector & kg \\
\(k_s\) & Modelled roughness height & m \\
\(C_s\) & Roughness constant & --- \\
\(\alpha_L\) & Bulk-liquid volume fraction & --- \\
\(\rho_L\) & Liquid density & kg/m³ \\
\(\boldsymbol u_L\) & Liquid velocity vector & m/s \\
\(\boldsymbol u_m\) & Mixture velocity vector & m/s \\
\(k\) & Shared turbulent kinetic energy & m²/s² \\
\(\epsilon\) & Shared turbulence dissipation rate & m²/s³ \\
\(S_L\) & Volumetric liquid mass source & kg/(m³·s) \\
\(V_c\) & Collector region / its geometric volume & m³ \\
\(\tau\) & Prescribed depletion time in the contact law & s \\
\(h_f\) & Film thickness & m \\
\(\tau_f\) & Direct film-drain capture time & s \\
\(K_d\) & Direct film-drain removal coefficient & s⁻¹ \\
\end{longtable}
}
\endgroup

\clearpage
\pagenumbering{arabic}

% Source: draft/01-introduction-and-literature-review.md
\section{Introduction}\label{introduction}

\subsection{Why separation matters}\label{why-separation-matters}

In a flash-steam geothermal power plant, hot fluid carries heat from an underground reservoir to the surface through production wells. As pressure falls, some liquid turns into steam, a process called flashing. Surface pipework carries the resulting steam--water mixture to a separator, where the flow divides into two routes: steam towards the power-generation equipment and liquid brine towards a separate outlet (\href{https://www.eia.gov/energyexplained/geothermal/geothermal-power-plants.php}{U.S. Energy Information Administration, n.d.}).

Figure 1 places the separator within this process, showing steam supplying the turbine--generator while brine and condensed exhaust steam return towards reinjection. The separator must keep these routes distinct: liquid entering the steam line can reach downstream equipment. The numbered circles identify thermodynamic states in the source; c1 and c2 indicate the condenser cooling-water connections. This is a plant schematic; the separator geometry used in this study is described below.

\par\nopagebreak[4]
\noindent\begin{minipage}{\linewidth}
\centering
\includegraphics[width=0.72\linewidth,height=0.344\textheight,keepaspectratio]{\detokenize{figures/literature/single-flash-plant-2017.png}}
\captionsetup{hypcap=false}\captionof{figure}[Single-flash geothermal plant]{Single-flash geothermal plant. Source: \href{https://doi.org/10.1186/s40517-017-0074-z}{El Haj Assad et al.~(2017)}, fig.~1, crediting Valdimarsson (2011).}
\label{fig:1}
\end{minipage}\par\medskip

This separation protects downstream equipment because liquid carried into the steam outlet can cause turbine erosion and transport minerals that form deposits. Separator design must therefore consider the quality of the outgoing steam together with pressure loss and operating pressure (\href{https://doi.org/10.1016/j.geothermics.2014.05.009}{Zarrouk \& Purnanto, 2015}).

This project concerns a vertical bottom-outlet cyclone (BOC) with a spiral inlet. The inlet directs the incoming steam--water mixture around a cylindrical vessel, creating a rotating flow. Within this flow, liquid moves towards the outer wall while steam flows towards an internal pipe that carries it out through the bottom. The term \emph{bottom outlet} refers to this steam-pipe arrangement; the brine leaves through a separate outlet (Figure 3). This arrangement determines how the rotating flow approaches the steam outlet (\href{https://doi.org/10.1016/j.geothermics.2014.05.009}{Zarrouk \& Purnanto, 2015}).

Transporting liquid towards the wall alone does not ensure effective separation. Wall liquid can form a film, circulate around the vessel, drain towards collection or become re-entrained into the steam as droplets. Stronger rotation can assist outward liquid transport, but steam shear can also promote droplet release. Improving a separator therefore requires understanding both liquid transport towards the wall and how it then drains or returns to the steam, rather than assuming that stronger swirl always gives better performance (Sections 2.2--2.4).

Computational fluid dynamics (CFD) can distinguish these competing liquid paths and help assess possible design changes. For that purpose, the model must represent both liquid capture and its later removal. This study examines how an inherited BOC model can be extended to meet that requirement.

\subsection{Starting model and its limits}\label{starting-model-and-its-limits}

The CFD domain represents the separator within the wider plant shown in Figure 1. Two simplifications interrupt the liquid path described above: the missing brine outlet and the lack of a separate wall-film model.

The first limitation is the lack of a separate liquid outlet: the inherited CFD geometry excludes the brine pool and its outlet pipe, leaving a closed lower wall. Without a way to remove liquid, incoming liquid must therefore build up within the model or leave through the steam outlet. Simply opening the lower boundary does not ensure that only liquid leaves, because steam can also escape through it. The simplified geometry lacks a route that removes collected liquid while retaining the intended steam path.

The second limitation is the unresolved behaviour of liquid after wall capture. The inherited formulation represents steam and liquid through a bulk Mixture field, without a separate model of an evolving wall film. Purnanto et al.~(2013), whose approach forms the basis of the inherited model, also tracked injected droplets using the Discrete Phase Model (DPM) to estimate outlet steam quality (Section 2.3.1). These trajectories describe droplet motion through the calculated flow, but do not model how deposited liquid is then stored, drained or released. Ordinary wall boundary conditions alone do not describe these film processes.

Together, these limitations leave the liquid path incomplete between wall capture and removal from the separator, so a low liquid flow through the steam outlet may reflect storage on the wall or at the base rather than sustained collection. The model must distinguish liquid transport, storage, release and removal before its predicted outlet flows can be used to assess separator performance.

\subsection{Aim and research questions}\label{aim-and-research-questions}

The aim is to develop the inherited separator CFD model and assess whether it provides a numerically consistent and physically credible representation of liquid transport and collection. The investigation addresses four questions:

\begin{enumerate}
\def\labelenumi{\arabic{enumi}.}
\item
  How do mesh representation and inlet phase placement affect the predicted liquid distribution, retention near the wall and drainage?
\item
  How can the inherited CFD model be improved to achieve greater numerical consistency and more realistic liquid behaviour?
\item
  How do wall roughness and wall-film modelling affect liquid transport, storage, release and collection?
\item
  How can CFD inform separator design optimisation to improve separation performance?
\end{enumerate}

These questions progress from understanding liquid distribution, through improving the model and examining wall transport, to identifying possible design changes. Progress is assessed using residuals, liquid mass balances, stored liquid mass (inventory) and liquid exit paths, while confidence in the predicted behaviour requires comparison with measurements under matching conditions and the same definitions of performance (Section 2.2). This study assesses the model improvements. The incomplete liquid mass balance and lack of field measurements under matching conditions limit claims about separator efficiency.

\subsection{Scope and contribution}\label{scope-and-contribution}

To address these questions, the study follows the liquid path through the simplified separator model. Mesh and inlet comparisons assess how liquid reaches and remains near the wall. Liquid-removal modelling provides the missing liquid exit, supported by startup procedures that allow the model to reach the intended feed conditions. Wall-roughness and wall-film studies then examine how captured liquid is stored, transported, released or drained. This understanding informs two inlet-region design proposals: a vertical film blocker and an angled film ramp. Initial tests examine local flow changes around these features, but have not shown improved separation. The simulations do not cover downstream equipment, mineral chemistry or the choice of optimum operating pressure.

Figure 2 brings these investigations together by distinguishing liquid transfer within the vessel from removal from the model. In this overview, liquid entering a wall film can remain in storage, drain away or return to the steam as droplets; these routes were examined in separate model configurations.

\par\nopagebreak[4]
\noindent\begin{minipage}{\linewidth}
\centering
\includegraphics[width=0.98\linewidth,height=0.43\textheight,keepaspectratio]{\detokenize{figures/report-liquid-destinations-overview.png}}
\captionsetup{hypcap=false}\captionof{figure}[Liquid destinations in the model]{Conceptual liquid pathways examined in separate model configurations.}
\label{fig:2}
\end{minipage}\par\medskip

The contribution is a set of checked liquid-removal and startup procedures, evidence of how model choices affect liquid transport, and two design proposals informed by the wall-film findings. Together, these extend the inherited model from predicting liquid distribution towards investigating how liquid reaches collection. The following literature review explains how earlier flow predictions, collection assumptions and film studies guide this development.

\section{Literature review}\label{literature-review}

\subsection{How the separator works}\label{how-the-separator-works}

A vertical bottom-outlet cyclone (BOC) uses a tangential or spiral inlet to rotate the incoming steam and brine, with liquid moving towards the wall while steam occupies the central region. An internal pipe carries steam out through the bottom; brine leaves separately. \emph{Bottom outlet} refers to the steam-pipe arrangement (\href{https://doi.org/10.1016/j.geothermics.2014.05.009}{Zarrouk \& Purnanto, 2015}).

Liquid at the wall can form a film that drains towards the brine outlet or releases droplets back into the steam (Figure 3). The diagram also distinguishes brine leaving the separator from liquid collected or condensed further downstream.

\par\nopagebreak[4]
\noindent\begin{minipage}{\linewidth}
\centering
\includegraphics[width=0.72\linewidth,height=0.344\textheight,keepaspectratio]{\detokenize{figures/literature/rizaldy-2016-figure-5-operation.png}}
\captionsetup{hypcap=false}\captionof{figure}[Separator liquid paths]{Conceptual schematic of separator liquid paths: wall film, carryover, brine discharge (\(\dot m_l\)) and downstream brine collection (\(\dot m_b\)). Source: \href{https://www.researchgate.net/publication/310818568_LIQUID_CARRYOVER_IN_GEOTHERMAL_STEAM-WATER_SEPARATORS}{Rizaldy et al.~(2016)}, fig.~5.}
\label{fig:3}
\end{minipage}\par\medskip

\subsection{Flow and performance}\label{flow-and-performance}

Rotation produces a low-pressure core, with inward acceleration scaling as \(u_\theta^2/r\), where \(u_\theta\) is the tangential velocity and \(r\) is the radius. The different responses of liquid and steam to pressure and drag, which depend on their densities, allow outward liquid transport. Because fine droplets can follow the steam whereas larger droplets reach collection surfaces more readily, uncertain inlet droplet sizes and continued droplet generation complicate prediction (\href{https://publications.mygeoenergynow.org/grc/1028587.pdf}{Pointon et al., 2009}).

Although higher velocity can strengthen outward transport, it can also increase entrainment by steam shear. This competing effect is reflected in empirical data showing sharp increases in outlet wetness at different inlet steam velocities for different separators (Figure 4). Thus, greater swirl alone does not establish better separation (\href{https://doi.org/10.1016/j.geothermics.2014.05.009}{Zarrouk \& Purnanto, 2015}).

\par\nopagebreak[4]
\noindent\begin{minipage}{\linewidth}
\centering
\includegraphics[width=0.72\linewidth,height=0.344\textheight,keepaspectratio]{\detokenize{figures/literature/zarrouk-purnanto-2015-figure-12-breakdown.png}}
\captionsetup{width=\linewidth}
\captionsetup{hypcap=false}\captionof{figure}[Inlet velocity and outlet wetness]{Outlet wetness and steam quality versus inlet steam velocity for three separator datasets. Left axis: outlet wetness; right axis: steam quality. Source: \href{https://doi.org/10.1016/j.geothermics.2014.05.009}{Zarrouk and Purnanto (2015)}, fig.~12, using Bangma (1960) and Lazalde-Crabtree (1984) data.}
\label{fig:4}
\end{minipage}\par\medskip

\par\noindent\begin{minipage}{\linewidth}
To compare these effects on the outgoing steam, outlet steam quality (the vapour mass fraction) is defined as

\[
x_{\mathrm{out}}=
\frac{\dot m_{v,\mathrm{out}}}
{\dot m_{v,\mathrm{out}}+\dot m_{l,\mathrm{out}}},
\tag{1}
\]

where the terms are outward vapour and liquid mass-flow rates through the steam outlet. Purnanto et al.~(2013) used this measure (\href{https://www.researchgate.net/publication/269519626_CFD_MODELLING_OF_TWO-PHASE_FLOW_INSIDE_GEOTHERMAL_STEAM-WATER_SEPARATORS}{Purnanto et al., 2013}). Outlet steam quality differs from the liquid-collection fraction, which compares collected liquid with inlet liquid; these measures use different reference mass flows and cannot be used interchangeably.
\end{minipage}\par\medskip

Steam quality also differs from steam purity, which concerns chemical contamination in the steam. Estimates of chemical carryover are affected by condensation and drain-pot collection, so field comparisons require matching definitions and sampling locations (\href{https://www.researchgate.net/publication/310818568_LIQUID_CARRYOVER_IN_GEOTHERMAL_STEAM-WATER_SEPARATORS}{Rizaldy et al., 2016}).

Assessing internal pressure drop requires defined inlet and outlet locations and a stated static or total pressure measure. This assessment differs from selecting the optimum operating pressure, which also involves flashing and plant conditions (\href{https://publications.mygeoenergynow.org/grc/1028587.pdf}{Pointon et al., 2009}; \href{https://doi.org/10.1016/j.geothermics.2014.05.009}{Zarrouk \& Purnanto, 2015}). The present model examines internal flow with specified steam and liquid conditions and does not simulate flashing.

\subsection{Previous design and CFD studies}\label{previous-design-and-cfd-studies}

Empirical geothermal design methods provide reference points for CFD, which can also reveal internal flow patterns. However, agreement with an empirical performance estimate does not establish that the model represents the correct liquid-collection mechanism (\href{https://doi.org/10.1016/j.geothermics.2014.05.009}{Zarrouk \& Purnanto, 2015}).

Pointon et al.~(2009) used Fluent's Discrete Phase Model (DPM) to compare inlet designs with assumptions about droplet feed and wall/baffle collection. Predicted efficiencies were 99.96\% for scrolled entry and 99.93\% for tangential entry; an empirical design method gave 99.955\%. Scrolled entry also gave a higher calculated mean tangential velocity than tangential entry: 34.1 rather than 32.5 m/s (Figure 5). However, agreement with the empirical method did not independently test the feed or collection assumptions (\href{https://publications.mygeoenergynow.org/grc/1028587.pdf}{Pointon et al., 2009}).

\par\nopagebreak[4]
\noindent\begin{minipage}{\linewidth}
\centering
\textbf{(a)}\par
\includegraphics[width=0.52\linewidth,height=0.18400000000000002\textheight,keepaspectratio]{\detokenize{figures/literature/pointon-2009-figure-5-inlets.png}}
\par\medskip\textbf{(b)}\par
\includegraphics[width=0.52\linewidth,height=0.18400000000000002\textheight,keepaspectratio]{\detokenize{figures/literature/pointon-2009-figure-7-velocity.png}}
\captionsetup{width=\linewidth}
\captionsetup{hypcap=false}\captionof{figure}[Inlet design and tangential velocity]{Scrolled entry gives a higher calculated mean tangential velocity than tangential entry. (a) Inlet geometries; (b) tangential-velocity fields. Scrolled entry (labelled ``lemniscate entrance'' in the source) is on the left; tangential entry is on the right. Source: \href{https://publications.mygeoenergynow.org/grc/1028587.pdf}{Pointon et al.~(2009)}, figs. 5 and 7.}
\label{fig:5}
\end{minipage}\par\medskip

Pointon et al.~(2009) and Purnanto et al.~(2013) both used renormalisation-group (RNG) \(k\)--\(\varepsilon\) turbulence modelling. Pointon preferred RNG with a swirl modification because it had lower computational requirements than Reynolds stress modelling. This lower computational cost helped guide the choice of turbulence model in this study, whose accuracy in swirling and near-wall flow requires further assessment (\href{https://publications.mygeoenergynow.org/grc/1028587.pdf}{Pointon et al., 2009}; \href{https://www.researchgate.net/publication/269519626_CFD_MODELLING_OF_TWO-PHASE_FLOW_INSIDE_GEOTHERMAL_STEAM-WATER_SEPARATORS}{Purnanto et al., 2013}).

Section 3 identifies the inherited modelling assumptions; Appendix A.1 records the controls used in the selected comparisons.

\subsubsection{Purnanto's model}\label{purnantos-model}

Purnanto et al.~compared three separator designs using steady, isothermal Mixture modelling, RNG \(k\)--\(\varepsilon\) turbulence and the SIMPLE pressure--velocity scheme. Walls were smooth; upstream pre-separation and brine discharge were excluded, with a constant pool level assumed (\href{https://www.researchgate.net/publication/269519626_CFD_MODELLING_OF_TWO-PHASE_FLOW_INSIDE_GEOTHERMAL_STEAM-WATER_SEPARATORS}{Purnanto et al., 2013, §§2.2--3.3, pp.~3--6}).

\par\noindent\begin{minipage}{\linewidth}
After the flow solution, DPM tracked nine Harwell-derived droplet sizes, uniformly distributed across the inlet. Although injection counts were equal, they represented unequal mass flows. For the collection estimate, droplets reaching the bottom were counted as separated, while outlet quality was steam mass divided by total outlet mass. The paper's earlier escaped/injected count definition instead describes penetration (\href{https://www.researchgate.net/publication/269519626_CFD_MODELLING_OF_TWO-PHASE_FLOW_INSIDE_GEOTHERMAL_STEAM-WATER_SEPARATORS}{Purnanto et al., 2013, §§2.5 and 4.3, pp.~4, 8}). The count-based quantity can be written as:

\[
P_N=\frac{N_{\mathrm{escaped}}}{N_{\mathrm{injected}}}.
\tag{2}
\]

Here, \(N_{\mathrm{escaped}}\) and \(N_{\mathrm{injected}}\) are particle counts, so this ratio differs from outlet steam quality, defined in Section 2.2, and from the fraction of inlet liquid collected. Droplet mass weighting is therefore essential when interpreting the reported performance.
\end{minipage}\par\medskip

The collection estimate also depended on incomplete tracks, which persisted even when the step limit rose from \(10^5\) to \(10^6\). Counting these incomplete tracks as separated improved empirical agreement, although the authors acknowledged that this assumption was not rigorous. Exact injection weights were not reported, and mesh refinement and experimental calibration remained necessary (\href{https://www.researchgate.net/publication/269519626_CFD_MODELLING_OF_TWO-PHASE_FLOW_INSIDE_GEOTHERMAL_STEAM-WATER_SEPARATORS}{Purnanto et al., 2013, pp.~8--9}).

Because vessel proportions also differed, the comparison could not isolate the effect of inlet shape. For the present study, a further limitation is the distinction between assigning droplets a destination and resolving liquid transport after wall capture. The approach did not model the wall film separately as it developed, so film storage, drainage and return to the steam could not be assessed. These limitations motivate explicit collection and wall-film treatments, while the Mixture and turbulence formulations provide a starting point for model development.

\subsection{Wall films and carryover}\label{wall-films-and-carryover}

After wall capture, film motion depends on gravity, pressure and shear, while its mass balance depends on deposition, drainage and entrainment. Deposition adds liquid to the film, whereas drainage removes it and entrainment returns droplets to the steam. More wall deposition can therefore increase liquid storage or later release without increasing external collection.

Rizaldy et al.~(2016) predict that entrainment is sensitive to loading, film thickness and inlet velocity. Figure 6 illustrates the conceptual mechanisms of undercut and roll-wave entrainment, associated with lower and higher film Reynolds numbers, respectively. Here, film Reynolds number compares inertia with viscosity using liquid velocity and film thickness. However, the brine-flow data needed to compare the predicted liquid carryover with measurements from the first drain pot after the Wairakei separator were unavailable, leaving quantitative agreement with field measurements unresolved (\href{https://www.researchgate.net/publication/310818568_LIQUID_CARRYOVER_IN_GEOTHERMAL_STEAM-WATER_SEPARATORS}{Rizaldy et al., 2016}).

\par\nopagebreak[4]\noindent\begin{minipage}{\linewidth}
\centering
\includegraphics[width=0.5760000000000001\linewidth,height=0.2752\textheight,keepaspectratio]{\detokenize{figures/literature/rizaldy-2016-figure-9-entrainment.png}}
\captionsetup{hypcap=false}\captionof{figure}[Droplet entrainment from a wall film]{Conceptual mechanisms by which steam shear releases droplets from a liquid film: (a) undercut and (b) roll-wave entrainment. Source: \href{https://www.researchgate.net/publication/310818568_LIQUID_CARRYOVER_IN_GEOTHERMAL_STEAM-WATER_SEPARATORS}{Rizaldy et al.~(2016)}, fig.~9, crediting Ishii and Grolmes (1975).}
\label{fig:6}
\end{minipage}\par\medskip

Fluent's roughness settings change near-wall momentum transfer; they do not add film storage or a drain (\href{https://ansyshelp.ansys.com/public/Views/Secured/corp/v252/en/flu_th/x1-5720008.164.html}{Ansys, Inc., 2025k}). Eulerian Wall Film (EWF) modelling, by contrast, adds surface storage and transport. These treatments answer different questions: roughness changes the interaction with the wall, while EWF represents the storage and transport of liquid on the wall. For this reason, the two treatments are tested separately using outlet flow, retained liquid mass and collection.

\subsection{Liquid flow and collection models}\label{liquid-flow-and-collection-models}

The Mixture model represents bulk liquid through the secondary-phase volume fraction, while algebraic slip estimates the relative velocity between phases under a local-equilibrium assumption. The liquid volume fraction does not directly resolve film thickness or the droplet-size distribution (\href{https://ansyshelp.ansys.com/public/Views/Secured/corp/v252/en/flu_th/flu_th_sec_drift_oview.html}{Ansys, Inc., 2025d}).

DPM instead tracks droplet paths that depend on droplet size. Interpreting these paths requires clear injection and boundary settings and a record of where each tracked droplet goes, including tracks that remain incomplete.

Skoog (2020) linked film and droplet behaviour through EWF surface transport, DPM deposition and entrainment in a reactor-channel approximation, comparing film and droplet flows with one-dimensional calculations using Okawa correlations based on experiments. The division of inlet liquid between film and droplets, together with the transverse droplet velocity, affected the result. For the present study, this provides a method for distinguishing film storage from droplet transport; using the reactor correlations for a cyclone still requires justification (\href{https://www.diva-portal.org/smash/get/diva2:1452100/FULLTEXT02.pdf}{Skoog, 2020}; \href{https://ansyshelp.ansys.com/public/views/secured/corp/v252/en/flu_ug/flu_ug_models_wallfilm.html}{Ansys, Inc., 2025e}).

When linking these models, internal transfer must be distinguished from removal: bulk-to-film transfer moves liquid within the domain, whereas the absorber removes it from the domain. Droplet injection must not duplicate liquid already supplied by another feed. Fluent permits cell-zone mass sources and associated transport sources (\href{https://ansyshelp.ansys.com/public/Views/Secured/corp/v252/en/flu_ug/flu_ug_bcs_sec_cell_zones.html}{Ansys, Inc., 2025a}), but the governing source law and its effects require assessment. Section 4.3 describes the contact-based absorber.

\subsection{Simulation setup and checks}\label{simulation-setup-and-checks}

In a steady calculation, iterations and pseudo-time describe numerical development, not physical filling time. The National Aeronautics and Space Administration (NASA) recommends monitoring residuals and relevant outputs, with convergence criteria based on acceptable error in those outputs (\href{https://www.grc.nasa.gov/www/wind/valid/tutorial/iterconv.html}{National Aeronautics and Space Administration {[}NASA{]}, 2021a}).

To prepare a starting flow field, Fluent permits flow development before film initialisation (\href{https://ansyshelp.ansys.com/public/views/secured/corp/v252/en/flu_ug/flu_ug_ewf_sec_overview.html}{Ansys, Inc., 2025h}). However, starting at reduced feed and gradually increasing it must be assessed for each case.

Implementation checks, iterative convergence, mass accounting and numerical sensitivity answer separate questions. In particular, mesh and timestep studies require a fixed model, whereas changes to roughness or the absorber test model sensitivity (\href{https://www.grc.nasa.gov/www/wind/valid/tutorial/verassess.html}{NASA, 2021c}).

Physical validation requires comparison with relevant measurements and assessment of uncertainty (\href{https://www.grc.nasa.gov/www/wind/valid/tutorial/valassess.html}{NASA, 2021b}). For this model, those comparisons must assess the assumed collection route and outlet predictions within the tested operating range.

\subsection{Research gaps}\label{research-gaps}

Pointon and Purnanto show how CFD can relate separator geometry to internal flow and predicted outlet performance (Section 2.3). The film studies add a further consideration: captured liquid can remain stored or return to the steam (Sections 2.4--2.5). Together, these approaches distinguish predicting wall capture from predicting sustained collection. Pointon's collecting boundaries and Purnanto's treatment of incomplete tracks also show why the assumed liquid destination matters when interpreting an outlet prediction.

The unresolved question for the inherited model is whether liquid delivered to the wall can reach collection without continued accumulation or return to the steam. The missing brine outlet makes a liquid-removal route necessary, while the film literature shows why liquid transport towards that route also needs study. This leads to the investigation in Section 4: compare liquid distribution, establish liquid removal and a practical startup procedure, examine wall-film transport, and use that understanding to propose design changes. Tracking storage and removal alongside outlet flow helps explain the full liquid path at each stage.

% Source: draft/02-03-model-development-methods-and-results.md
\section{Starting model and setup}\label{starting-model-and-setup}

The literature review identified a gap between predicting liquid capture at the wall and predicting sustained collection (Section 2.7). To address that gap, this study extended the inherited Purnanto-based model, using its geometry and bulk-flow formulation as a starting point. The missing brine outlet and lack of a separate wall-film model were the main areas for development.

Section 3 defines the geometry, bulk-flow model, operating conditions and approach to comparisons. Section 4 describes the model development; Appendix A gives detailed settings.

\subsection{Geometry and model scope}\label{geometry-and-model-scope}

\par\noindent\begin{minipage}{\linewidth}
Based on Purnanto et al.~(2013), the inherited geometry retained the spiral inlet, cylindrical separation vessel and internal steam outlet pipe needed to study liquid transport through the active separation region (Figure 7). In this simplified domain, the lower brine pool and its outlet pipe were excluded, leaving a closed bottom wall.

\medskip
\centering
\includegraphics[width=0.98\linewidth,height=0.43\textheight,keepaspectratio]{\detokenize{figures/model-domain-overview.png}}
\captionsetup{hypcap=false}\captionof{figure}[Computational domain and inlet arrangement]{Simplified computational domain and inlet arrangement: (a) spiral inlet, vessel and internal steam outlet; (b) wall-side liquid strip and adjacent vapour inlet. The brine pool and discharge pipe are omitted. Panel (b) shows the split inlet introduced in Section 4.1; the absorber is introduced in Section 4.3.}
\label{fig:7}
\end{minipage}\par\medskip

The split inlet shown in Figure 7(b), like the absorber, was introduced during development rather than inherited from the original mixed-feed arrangement. The closed base provided no physical liquid outlet. In the later liquid-removal studies, the absorber removed liquid from a defined lower region while the base kept steam on its intended path (Section 4.3).

\subsection{Bulk-flow model and assumptions}\label{bulk-flow-model-and-assumptions}

The Mixture model carried forward Purnanto's modelling approach. Its algebraic-slip formulation followed that study's small-Stokes-number argument: sufficiently small dispersed droplets respond rapidly relative to the characteristic flow timescale (\href{https://www.researchgate.net/publication/269519626_CFD_MODELLING_OF_TWO-PHASE_FLOW_INSIDE_GEOTHERMAL_STEAM-WATER_SEPARATORS}{Purnanto et al., 2013}). Initial volume of fluid (VOF) and turbulence trials did not establish Mixture with RNG \(k\)--\(\varepsilon\) as the most accurate formulation. It was retained because solving mixture momentum with algebraic slip requires fewer equations than a full Eulerian formulation, making repeated inlet and wall-treatment studies practical (\href{https://ansyshelp.ansys.com/public/Views/Secured/corp/v252/en/flu_th/flu_th_sec_drift_oview.html}{Ansys, Inc., 2025d}).

A pressure-based solver with the Mixture model in Ansys Fluent 2025 R2 was used for bulk continuity, momentum, liquid fraction and algebraic slip, with steam as the primary phase and water as the secondary phase. The inherited setup carried forward Purnanto's SIMPLE pressure--velocity scheme, while Appendix A.1 identifies comparison cases that used the Coupled pressure--velocity scheme. Turbulence was represented by the RNG \(k\)--\(\varepsilon\) model with swirl treatment and standard wall functions, and gravity acted downwards. With the energy equation disabled, flashing, condensation and temperature-dependent properties were excluded. Table 1 gives the material values, while Appendix A.1 records the controls (\href{https://ansyshelp.ansys.com/public/Views/Secured/corp/v252/en/flu_th/flu_th_sec_multiphase_mixture.html}{Ansys, Inc., 2025d}).

\begin{table}[!htbp]
\centering
\caption{Fixed material properties used in the throughput-controlled absorber model.}\label{tab:1}
\begingroup\small\setlength{\tabcolsep}{3pt}
{
\begin{tabular}{@{}
  >{\raggedright\arraybackslash}p{(\linewidth - 4\tabcolsep) * \real{0.2500}}
  >{\raggedleft\arraybackslash}p{(\linewidth - 4\tabcolsep) * \real{0.3200}}
  >{\raggedleft\arraybackslash}p{(\linewidth - 4\tabcolsep) * \real{0.4300}}@{}}
\toprule\noalign{}
\begin{minipage}[b]{\linewidth}\raggedright
Phase
\end{minipage} & \begin{minipage}[b]{\linewidth}\raggedleft
Density (kg/m³)
\end{minipage} & \begin{minipage}[b]{\linewidth}\raggedleft
Dynamic viscosity (Pa·s)
\end{minipage} \\
\midrule\noalign{}
Vapour & 5.79743 & \(1.52062\times10^{-5}\) \\
Liquid water & 881.211 & \(1.45544\times10^{-4}\) \\
\bottomrule\noalign{}
\end{tabular}
}
\endgroup
\end{table}

The bulk-flow formulation used Manninen slip, Schiller--Naumann drag and a liquid-phase diameter of 10 µm (Appendix A.1).

\subsection{Inlet and boundary conditions}\label{inlet-and-boundary-conditions}

The model comparisons used Purnanto's reference feed condition: 197.61 kg/s at 1600 kJ/kg and 11.2 bar absolute. At this pressure, saturated liquid and vapour enthalpies, \(h_f\) and \(h_g\), gave the vapour mass fraction as \(x=(h-h_f)/(h_g-h_f)=0.408323\). The inlet boundary conditions used Purnanto's rounded phase mass-flow rates of 116.92 kg/s liquid and 80.69 kg/s vapour; Appendix A.2 gives the calculation (\href{https://www.researchgate.net/publication/269519626_CFD_MODELLING_OF_TWO-PHASE_FLOW_INSIDE_GEOTHERMAL_STEAM-WATER_SEPARATORS}{Purnanto et al., 2013}). Flashing was not solved within the domain. Enthalpy was therefore used to set the inlet steam--water composition for the model-development comparisons.

Stationary, no-slip walls confined the flow, and the steam outlet supplied the pressure boundary condition, so inlet pressure was calculated by the solver. With the operating pressure set to 0 Pa, outlet gauge and absolute pressures were equal. The reference feed condition defines full feed in the following sections; reduced-feed holds and ramps changed the inlet flow during startup, while the inlet comparison changed where the phases entered. Appendix A.1 gives the turbulence and backflow inputs and the differences between comparison setups.

\subsection{Study plan and comparisons}\label{study-plan-and-comparisons}

The investigation progressed from liquid distribution and removal to wall-film transport, mesh sensitivity and proposed designs. The 60k poly-hexcore mesh supported repeated model development, followed by tests on finer meshes; the historical tetrahedral case provided context. Comparisons used shared starting fields and comparison intervals where possible, with preparation differences identified separately. Table A1a in Appendix A.1 summarises the meshes, settings and comparison limits.

Scaled continuity residuals, liquid inventories, removal rates and contours were assessed together to interpret each test. In the following sections, outlet liquid flow means the liquid mass-flow rate through the steam outlet and is positive outwards unless stated otherwise. Bulk-flow development is reported in solver iterations; film development is reported in accepted film time, the simulated time advanced by the film solver.

\section{Model development and results}\label{model-development-and-results}

The results address the four research questions through inlet and mesh comparisons, numerical improvements, wall-liquid transport studies and initial design tests. Table A1a summarises the comparisons.

\subsection{Split inlet and liquid distribution}\label{split-inlet-and-liquid-distribution}

The split inlet was introduced to represent the expected wall-side concentration of liquid as the upstream flow turned into the real separator. Using the geometry in Section 3.1 (Figure 7), the 0.724 m × 0.724 m inlet cross-section was divided into a 6.754 mm wall-side liquid strip and the remaining steam inlet. The reference feed defined in Section 3.3 supplied approximately 116.9 kg/s of liquid and 80.7 kg/s of steam through the liquid and steam inlet faces, respectively. Complete phase separation at entry was an idealisation of this expected distribution; Appendix A.2 explains how the strip width was chosen. With this physical intent represented in the inlet, the investigation then examined how the change affected liquid behaviour within the vessel.

Matched cases compared mixed and split inlets under the controls in Table A1a. Although boundary mass-flow rates changed little relative to mixed feed, with almost all incoming liquid still leaving through the steam outlet, the contours showed a clearer spatial difference. The split-inlet case has a more continuous liquid-rich band around the outer wall in Figure 8(b), while the vertical section in Figure 8(d) shows liquid concentrated beside the wall and in a deeper region near the base. The central bulk flow retains a low liquid volume fraction.

\par\nopagebreak[4]
\noindent\begin{minipage}{\linewidth}
\centering
\includegraphics[width=0.98\linewidth,height=0.43\textheight,keepaspectratio]{\detokenize{figures/results-inlet-spatial-comparison.png}}
\captionsetup{hypcap=false}\captionof{figure}[Liquid distribution with mixed and split feed]{Liquid volume fraction with mixed (a, c) and split inlets (b, d), at inlet height (a, b) and through the vessel centre (c, d).}
\label{fig:8}
\end{minipage}\par\medskip

The closed base still had no liquid outlet, as expected from the simplified geometry. The contribution of the split inlet was to represent the intended wall-side liquid entry and show how it changed the liquid distribution inside the separator. The clearer wall band and deeper base accumulation gave a spatial basis for examining liquid retention on different meshes.

\subsection{Liquid accumulation on different meshes}\label{liquid-accumulation-on-different-meshes}

The mesh investigation moved from the historical tetrahedral mesh to poly-hexcore meshes to improve the representation of swirling flow and liquid transport near the wall. The comparisons examined how this change affected liquid distribution and retention.

Poly-hexcore concentrated resolution near the wall through wall-following prism layers, connected by polyhedral cells to larger hexahedral cells in the core. This allowed small wall-normal spacing without requiring the same resolution throughout the vessel (\href{https://www.ansys.com/content/dam/resource-center/white-paper/ansys-fluent-mosaic-technology-wp.pdf}{Ansys, Inc., 2020, pp.~1--2}).

\FloatBarrier
\par\noindent\begin{minipage}{\linewidth}
Figure 9 compares mesh sections and liquid fields from the historical tetrahedral case and a 60k poly-hexcore case at iteration 10,000. The historical case shows a thin liquid-rich region near the wall. The poly-hexcore case shows a wider region with a higher liquid fraction, revealing a clear difference in wall-liquid retention.

\medskip
{\centering
\includegraphics[width=0.98\linewidth,height=0.31\textheight,keepaspectratio]{\detokenize{figures/results-mesh-spatial-comparison.png}}
\par}
\captionsetup{hypcap=false}\captionof{figure}[Historical tetrahedral and poly-hexcore mesh sections (a--b) and liquid-volume-fraction fields (c--d).]{Historical tetrahedral and poly-hexcore mesh sections (a--b) and liquid-volume-fraction fields (c--d).}
\label{fig:9}
\end{minipage}\par\medskip


Separate split-inlet cases on the 342k, 680k and 997k meshes showed how retained liquid developed during calculation. Liquid accumulated near the wall and base, and its inventory continued to rise while little escaped through the steam outlet (Figure 10). With no separate liquid exit at the closed base, these cases retained liquid instead of reaching a steady inventory.

\par\nopagebreak[4]
\noindent\begin{minipage}{\linewidth}
\centering
\includegraphics[width=0.88\linewidth,height=0.43\textheight,keepaspectratio]{\detokenize{../../experiments/phase-08-storyline-reconstruction/stage-02-fine-mesh-simple/direct-student-342k/n3000/figures/matched-mesh-checkpoints.png}}
\captionsetup{hypcap=false}\captionof{figure}[Liquid accumulation on finer meshes]{Liquid inventory (left) and outlet liquid flow (right) on the 342k, 680k and 997k meshes.}
\label{fig:10}
\end{minipage}\par\medskip

Continued calculation on the 680k mesh produced further accumulation before numerical divergence stopped the run, whereas the historical tetrahedral case had continued much further. Seen alongside the wider wall-liquid region in Figure 9, this contrast raised the possibility that the poly-hexcore meshes retained liquid that the earlier mesh had spread into the bulk flow.

Fluent's guidance links reduced numerical spreading to cell refinement and alignment with the flow. Wall-following layers do not automatically align with the spiralling liquid path, and tetrahedral meshes can also contain prism layers; cell spacing and orientation therefore matter more than the mesh label alone (\href{https://ansyshelp.ansys.com/public/Views/Secured/corp/v252/en/flu_ug/flu_ug_GridTypes.html}{Ansys, Inc., 2025c, §7.1.3.3}). Core resolution also matters because swirl and pressure gradients drive the surrounding flow (\href{https://ansyshelp.ansys.com/public/Views/Secured/corp/v252/en/flu_ug/flu_ug_uns_sec_models_swirl.html}{Ansys, Inc., 2025l, §§10.3.2--10.3.3}).

The comparisons do not isolate mesh construction from all other differences, so reduced spreading remains a possible explanation. What they did reveal was sustained wall and base accumulation that the earlier calculation had not shown as clearly. This gave a clearer view of where liquid remained within the separator and a basis for studying its removal in Section 4.3.

\subsection{Adding liquid removal}\label{adding-liquid-removal}

The next development added a destination for liquid retained near the base. Earlier trials included the lower geometry to model the brine pool and its outlet directly. Across the two sets of steady outlet tests, only six of 16 cases reached 500 iterations, and completed cases retained large liquid mass-balance errors (Appendix A.7). These difficulties led to a simpler test: provide liquid removal while studying transport separately from pool behaviour and discharge control.

The absorber removed liquid reaching the lower region, while the closed bottom diverted steam towards its outlet. The rigid, no-slip bottom remained an approximation to the real steam--brine interface, and pool motion and control of liquid outflow were not modelled. This arrangement supplied a liquid destination near the base while preserving the steam's turn towards its outlet, allowing liquid transport towards the removal region to be tested.

\par\noindent\begin{minipage}{\linewidth}
The absorber removed liquid through a mass sink in a defined cell region near the base. Its first form targeted the incoming liquid throughput and was used for the startup in Section 4.4. The selected contact-based absorber instead made removal depend on liquid reaching the absorber, using local liquid content and a specified depletion time:

\[
S_L=-\frac{\rho_L\alpha_L}{\tau},\qquad \tau=10^{-5}\ \mathrm{s}.
\tag{3}
\]

Here, \(\rho_L\) is liquid density, \(\alpha_L\) is liquid volume fraction and \(\tau\) is the depletion time; \(S_L\) is the mass sink per unit volume, in kg/(m³·s). The depletion time set the rate of liquid removal, and tests of removal strength guided its selection. Associated momentum and turbulence sources represented liquid removal without a direct steam mass sink. Appendix A.3 gives the source laws, mass-balance checks and depletion-time test results (\href{https://ansyshelp.ansys.com/public/Views/Secured/corp/v252/en/flu_ug/flu_ug_bcs_sec_cell_zones.html}{Ansys, Inc., 2025a}).
\end{minipage}\par\medskip

\subsection{Starting at reduced feed}\label{starting-at-reduced-feed}

The inlet was held at 25\% of the reference feed with liquid removal active, allowing the continuity residual to fall before more model features were added. This earlier startup used the 60k mesh, Hybrid Initialisation, smooth walls and no wall-film model, together with a throughput-controlled absorber. In this absorber, total removal targeted the inlet liquid mass flow and was distributed according to liquid content in the absorber region. Table 2 gives the reduced and full-feed boundary conditions.

\begin{table}[!htbp]
\centering
\caption{Main boundary conditions. Reduced feed was 25\% of full feed, with the phase mass ratio retained.}\label{tab:2}
\begingroup\small\setlength{\tabcolsep}{3pt}
{
\begin{tabular}{@{}
  >{\raggedright\arraybackslash}p{(\linewidth - 6\tabcolsep) * \real{0.2000}}
  >{\raggedright\arraybackslash}p{(\linewidth - 6\tabcolsep) * \real{0.2900}}
  >{\raggedleft\arraybackslash}p{(\linewidth - 6\tabcolsep) * \real{0.2900}}
  >{\raggedleft\arraybackslash}p{(\linewidth - 6\tabcolsep) * \real{0.2200}}@{}}
\toprule\noalign{}
\begin{minipage}[b]{\linewidth}\raggedright
Boundary
\end{minipage} & \begin{minipage}[b]{\linewidth}\raggedright
Condition
\end{minipage} & \begin{minipage}[b]{\linewidth}\raggedleft
Full feed
\end{minipage} & \begin{minipage}[b]{\linewidth}\raggedleft
Reduced feed
\end{minipage} \\
\midrule\noalign{}
Liquid inlet & Liquid mass-flow inlet & 116.9200 kg/s & 29.2300 kg/s \\
Vapour inlet & Vapour mass-flow inlet & 80.6900 kg/s & 20.1725 kg/s \\
Steam outlet & Gauge pressure & 1.120 MPa & Same \\
Vessel and bottom & Stationary, no-slip walls & Selected roughness/film treatment & As defined for startup \\
\bottomrule\noalign{}
\end{tabular}
}
\endgroup
\end{table}

The reduced-feed hold produced a more stable liquid inventory. The estimated inventory fall was only 0.30\% between iterations 1080 and 1580, with about 31 kg remaining at the end of the hold. This was approximately 93\% below the 454 kg stored in the absorber-free 342k reference, whose inventory was still rising. The two cases differed in mesh and loading, so this contrast reflects the combined changes in removal, feed and mesh (Appendix A.3).

The scaled continuity residual also fell, reaching a minimum of \(3.2\times10^{-4}\) (Figure 11), about 650 times lower than the approximately 0.21 at the absorber-free 342k comparison point (Appendix A.3). The lower reported residual and small inventory change together supported gradually returning to full feed. Each loading stage started from the flow field reached in the previous stage. This gradual increase is consistent with Fluent guidance for improving solution stability in swirling flows (\href{https://ansyshelp.ansys.com/public/Views/Secured/corp/v252/en/flu_ug/flu_ug_uns_sec_models_swirl.html}{Ansys, Inc., 2025l}).

\par\nopagebreak[4]
\noindent\begin{minipage}{\linewidth}
\centering
\includegraphics[width=0.98\linewidth,height=0.43\textheight,keepaspectratio]{\detokenize{figures/results-described-early-development.png}}
\captionsetup{hypcap=false}\captionof{figure}[Reduced-feed startup and development]{Residuals, liquid inventory and inlet loading during the reduced-feed hold, ramp and subsequent development.}
\label{fig:11}
\end{minipage}\par\medskip

Figure 11 shows the continuity residual rising during the ramp, then falling to about \(2.8\times10^{-3}\) during full-feed development with Coupled pressure--velocity scheme and Global Time Step pseudo-time. Markers A, B and C identify the ramp start, full feed with Coupled activation, and wall-film activation, respectively; the horizontal axis records solver iterations, not physical time.

Residual oscillations and an unresolved liquid mass balance remained, so the full-feed state was not treated as converged. Nevertheless, the startup showed substantially improved numerical behaviour: the continuity residual decreased markedly, inventory changed little late in the reduced-feed hold, and the calculation reached full feed. These improvements justified saving the smooth-wall, film-disabled state at iteration 5586 as the common baseline for the wall-treatment comparisons.

\subsection{Liquid remaining near the upper wall}\label{liquid-remaining-near-the-upper-wall}

The absorber provided a lower liquid-removal route, but substantial liquid still reached the steam outlet at full feed. Figure 12 illustrates the concern in an earlier absorber trial: liquid-rich bands extend up the outer walls to the upper shoulders, with smaller liquid-rich regions near the internal outlet entrance.

\par\nopagebreak[4]
\noindent\begin{minipage}{\linewidth}
\centering
\includegraphics[width=0.98\linewidth,height=0.43\textheight,keepaspectratio]{\detokenize{../../meetings/2026-09-14/figures/P07A-02-absorber-active1900-XY-VOF-velocity.png}}
\captionsetup{hypcap=false}\captionof{figure}[Wall liquid in an earlier absorber trial]{Liquid volume fraction and velocity vectors in an earlier absorber trial.}
\label{fig:12}
\end{minipage}\par\medskip

Providing a sink addressed local removal, but effective collection also required liquid transport towards that region. The upper-wall liquid suggested that gravity and wall friction were not strong enough to redirect it downwards. Although the contours alone did not establish the force balance, they motivated the roughness and wall-film studies in Sections 4.6.1 and 4.6.2, which examined whether changing the wall treatment could help liquid reach the absorber.

\subsection{Wall--liquid interaction}\label{wallliquid-interaction}

Two separate tests from the common baseline examined the wall-liquid problem: one varied roughness to test whether changing wall friction altered the liquid path, while the other introduced EWF to represent liquid as a film with its own storage and transport. To separate these treatments before examining combined cases, the roughness tests kept EWF disabled and the film tests retained smooth walls.

Each test restarted independently from iteration 5586 and ran for 3000 updates. Means were calculated over the final 500 updates, keeping the starting field and comparison interval the same (Section 3.4). Inventory was assessed alongside outlet flow to distinguish changed routing from depletion of stored liquid.

\subsubsection{Wall roughness}\label{wall-roughness}

Wall roughness was included in the wall-friction treatment through the equivalent sand-grain height \(k_s\) and roughness constant \(C_s\) (\href{https://ansyshelp.ansys.com/public/Views/Secured/corp/v252/en/flu_th/x1-5720008.164.html}{Ansys, Inc., 2025k}). The sensitivity tests examined how these parameters changed liquid flow through the steam outlet, the primary quantity used here to investigate carryover. Appendix A.6 records the wall treatment.

Roughness ranged from 0 to 8 mm, with detailed comparison controls in Appendix A.1.

Changing wall roughness produced large changes in mean bulk-liquid flow through the steam outlet (Figure 13). At \(C_s=0.5\), a roughness height of 0.05 mm raised mean outlet flow by about 21\% relative to the smooth-wall reference, whereas 8 mm reduced it by about 44\%. Increasing \(C_s\) lowered the flow at both fixed roughness heights. Dashed lines mark the smooth-wall reference.

\par\nopagebreak[4]
\noindent\begin{minipage}{\linewidth}
\centering
\includegraphics[width=0.98\linewidth,height=0.43\textheight,keepaspectratio]{\detokenize{figures/results-described-roughness.png}}
\captionsetup{hypcap=false}\captionof{figure}[Wall roughness and outlet liquid flow]{Mean outlet liquid flow versus (a) roughness height and (b) roughness constant.}
\label{fig:13}
\end{minipage}\par\medskip

The 8 mm case also reduced stored liquid mass substantially, showing that wall friction affected both storage and outlet flow. Because increased roughness could raise or lower outlet flow, its selected value needed a physical basis. A commercial-steel reference, \(k_s=0.045\) mm and \(C_s=0.5\), was selected to represent the intended wall material, with the bottom kept smooth (Appendix A.1).

\subsubsection{Adding wall-film modelling}\label{adding-wall-film-modelling}

Changing wall friction did not explicitly represent liquid stored on the wall. The second test therefore returned to the smooth-wall baseline and introduced the Eulerian Wall Film (EWF) model. It solved surface mass and momentum equations, with accretion transferring bulk liquid into the film (\href{https://ansyshelp.ansys.com/public/Views/Secured/corp/v252/en/flu_ug/flu_ug_ewf_sec_options.html}{Ansys, Inc., 2025e}).

With common smooth-wall starting fields and throughput-controlled liquid removal, EWF with accretion reduced mean bulk-liquid outlet flow by about 93\%, from 24.4 to 1.74 kg/s over the matched reporting window (Figure 14). Film mass and momentum were solved together at a 10 µs step. The shaded region marks the final 500 iterations used for comparison.

\par\nopagebreak[4]
\noindent\begin{minipage}{\linewidth}
\centering
\includegraphics[width=0.784\linewidth,height=0.34400000000000003\textheight,keepaspectratio]{\detokenize{figures/results-described-film-outlet.png}}
\captionsetup{hypcap=false}\captionof{figure}[Outlet flow with film disabled and enabled]{Outlet liquid flow with EWF disabled and accretion enabled.}
\label{fig:14}
\end{minipage}\par\medskip

During the accretion-enabled continuation, bulk liquid mass fell from about 296 to 63 kg, while film mass reached about 3 kg. The marked reduction in steam-outlet liquid flow showed how strongly wall-film transfer affected predicted carryover. With liquid now entering a separate film, longer calculations were needed to establish where it went: spreading along the wall, remaining stored or draining.

Combining roughness, film transport and the contact-based absorber still left drainage about 10\% below accretion over the final window (Appendix A.1). This continued storage motivated longer film studies; first, the lengthy startup was revised to introduce these treatments earlier.

\subsection{Shorter startup with wall-film modelling}\label{shorter-startup-with-wall-film-modelling}

Before extending film development, the long startup procedure was revised to introduce the combined treatments earlier and reach full feed with smaller loading disturbances. The shortened procedure reused the reduced-feed field without reinitialisation, then activated the contact-based absorber, roughness, the Coupled scheme and an initially dry film before increasing the feed.

The original and shortened startup sequences were compared as complete procedures (Section 3.4 and Appendix A.1).

The shortened startup substantially reduced disturbances during the loading ramp. The continuity normalisation was identical for both startup ramps (Appendix A.3). Within the ramp window, the peak continuity residual fell by 70.15\%, peak combined bulk and film inventory by 70.43\%, and peak outward liquid flow through the steam outlet by 75.78\% relative to the historical ramp (Figure 15).

\par\nopagebreak[4]
\noindent\begin{minipage}{\linewidth}
\centering
\includegraphics[width=0.98\linewidth,height=0.43\textheight,keepaspectratio]{\detokenize{figures/results-described-startup.png}}
\captionsetup{hypcap=false}\captionof{figure}[Historical and shortened startup ramps]{Historical and shortened startup: (a) continuity residual, (b) liquid inventory and (c) outlet liquid flow.}
\label{fig:15}
\end{minipage}\par\medskip

Switching on the added treatments during the activation hold produced a continuity residual peak of 0.10593, higher than the historical ramp peak. This hold is excluded from Figure 15. The film solution did not meet the thickness and momentum tolerance in 15 of the final 100 hold updates, but met it throughout the following ramp (Appendix A.4). The shortened procedure thus separated the difficult activation step from the feed increase, producing a smoother loading ramp.

The check of the contact-based absorber also confirmed that removal followed the specified law (Appendix A.3). By activating the treatments before increasing feed, the revised procedure produced a smoother transition to full loading, with lower peaks in residual, liquid inventory and carryover during the ramp. This improved startup supplied the full-feed starting field for the longer film studies.

\subsection{Film growth and drainage}\label{film-growth-and-drainage}

The initial film produced during startup was advanced to 500 ms to examine whether it would spread and develop sustained drainage. The surrounding bulk flow field was held fixed so that its effect on the film remained unchanged. This bulk field had not converged beforehand. Drainage here means Fluent's reported liquid outflow across the film-domain edges, excluding the bulk absorber and the direct film drain introduced in Section 4.9 (Appendix A.5).

Drainage rose rapidly during early development and then began to level out (Figure 16), but remained below film input (accretion) as inventory continued to rise through 500 ms. Over the final reporting window, drainage was still approximately 14\% below input. The film mass-balance check supported this accumulation: over 490--500 ms, input minus drainage agreed with the measured storage rate to within 0.008\% of input (Appendix A.5). The nearly level drainage curve therefore hid continued storage, which the mass balance and inventory history revealed together.

\par\nopagebreak[4]
\noindent\begin{minipage}{\linewidth}
\centering
\includegraphics[width=0.98\linewidth,height=0.43\textheight,keepaspectratio]{\detokenize{figures/results-film-development-reference500ms.png}}
\captionsetup{hypcap=false}\captionof{figure}[Film storage and drainage to 500 ms]{Film inventory, input and drainage through 500 ms.}
\label{fig:16}
\end{minipage}\par\medskip

The contour plots show film development across the vessel wall (Figure 17). Spreading and thickening in the lower and middle wall regions show that film growth involved both increasing surface coverage and increasing local storage.

\par\nopagebreak[4]
\noindent\begin{minipage}{\linewidth}
\centering
\includegraphics[width=0.98\linewidth,height=0.43\textheight,keepaspectratio]{\detokenize{figures/wall-film-described-development.png}}
\captionsetup{hypcap=false}\captionof{figure}[Wall-film development under frozen bulk]{Wall-film thickness during development with the bulk flow held fixed.}
\label{fig:17}
\end{minipage}\par\medskip

Small film timesteps made these calculations costly. In a matched 2.5 ms test, 20 µs passed the local checks, while 25 µs exceeded the 0.1\% film mass-balance error limit despite similar film fields. The mass-balance check therefore determined which timestep passed the test; longer runs still required checks on stability and mass balance as the film developed (Appendix A.4).

Together, the histories and contours showed how liquid spread, built up and drained along the wall. Tracking film mass alongside outflow revealed continued storage even after drainage began to level out. This linked film development to the balance between supply and removal, directing the next tests towards release and direct drainage.

\subsection{Film release and lower-wall drainage}\label{film-release-and-lower-wall-drainage}

Film liquid continued to accumulate, so two ways to reduce it were tested. First, stripping and edge separation allowed liquid to leave the film as droplets (\href{https://ansyshelp.ansys.com/public/Views/Secured/corp/v252/en/flu_th/flu_th_ewf_sec_film_submodels.html}{Ansys, Inc., 2025b}). Film mass still increased (Appendix A.5). The next test added a drain at the lower wall.

Because the bulk absorber could not remove film liquid, the film domain was extended onto the lower wall and a separate mass and momentum sink was added as a direct film drain. Matched Drain OFF/ON cases used identical starting fields, the same film domain and bulk flow held fixed. Over 15 ms, the drain reduced lower-wall inventory by about 70\% (Figure 18). Net extra external removal was 0.1897 kg, closely matching the reduction in total film inventory and confirming the added liquid removal over this interval (Appendix A.5).

\par\nopagebreak[4]
\noindent\begin{minipage}{\linewidth}
\centering
\includegraphics[width=0.98\linewidth,height=0.43\textheight,keepaspectratio]{\detokenize{figures/results-direct-film-drain.png}}
\captionsetup{hypcap=false}\captionof{figure}[Direct drainage and external removal]{Drain OFF/ON comparison: (a) lower-wall film mass and (b) external liquid removal.}
\label{fig:18}
\end{minipage}\par\medskip

During longer development, lower-wall film inventory remained low while upper-film inventory continued to grow. The later Courant-limit failure means that the end of the run cannot be treated as a reliable endpoint (Appendix A.5). Upper-wall accumulation was already present before that failure. The inventory trends show that the remaining storage problem was above the drain: removal controlled the lower region while liquid continued to accumulate higher on the wall.

The drain removed liquid that reached the lower wall. The remaining task was to understand how liquid reached it from higher in the vessel. This result separated the two parts of collection---transport to the drain and removal at the drain---and gave the next mesh comparison a clear focus.

\subsection{Numerical checks and mesh sensitivity}\label{numerical-checks-and-mesh-sensitivity}

The coarse mesh provided a practical basis for developing film transport and drainage. The next question was whether its predicted film coverage, storage and drainage remained similar on finer meshes. Four meshes (60k, 342k, 680k and 997k) were compared at a common film time of 195 ms. Earlier coarse-mesh drainage had begun to level out around 200 ms (Figure 16), providing a practical observation point for comparing film paths, storage and removal across meshes.

Each film evolved with the bulk flow for its mesh held fixed and the selected steel roughness applied. The prepared bulk flows and starting film states differed between meshes (Appendix A.6), so the films began under different conditions. The comparison examined how each prepared flow supported film spreading, storage and drainage.

Figure 19 unfolds the cylindrical wall, with height vertically and angle around the vessel horizontally. The white opening marks the inlet. All panels use a 0--0.30 mm thickness scale, with higher values clipped; purple indicates low thickness, not necessarily a dry wall. Arrows show direction rather than speed.

\par\nopagebreak[4]
\noindent\begin{minipage}{\linewidth}
\centering
\includegraphics[width=0.98\linewidth,height=0.50\textheight,keepaspectratio]{\detokenize{../../experiments/phase-09-mesh-convergence/figures/selected-film-thickness.png}}
\captionsetup{hypcap=false}\captionof{figure}[Wall-film distribution across meshes]{Wall-film thickness before development (top) and at 195 ms (bottom).}
\label{fig:19}
\end{minipage}\par\medskip

At 195 ms, relatively thick film covers much of the displayed wall on the 60k mesh. The finer meshes show more distinct liquid paths and thinner upper-wall regions. The 342k and 680k patterns are similar, while the 997k case has less clearly visible film above the inlet. These paths show where liquid travels along the wall, making their position and development relevant to drainage.

The inventory and drainage results support this spatial contrast (Figure 20). The finer meshes contain less film and have lower drain rates than the 60k case. Although 342k and 680k give similar film inventories and drainage, the 997k case contains about 30\% less film and drains about 32\% less than 680k. The finest completed mesh breaks the apparent agreement between the intermediate meshes, so mesh independence has not been established. The fall in both film mass and drainage on the finest mesh links the changed removal rate to a different wall-liquid inventory.

\par\nopagebreak[4]
\noindent\begin{minipage}{\linewidth}
\centering
\includegraphics[width=0.98\linewidth,height=0.43\textheight,keepaspectratio]{\detokenize{../../experiments/phase-09-mesh-convergence/figures/selected-mesh-response.png}}
\captionsetup{hypcap=false}\captionof{figure}[Liquid inventories and removal across meshes]{Bulk-liquid quantities, film inventory and lower-wall drain rate across meshes.}
\label{fig:20}
\end{minipage}\par\medskip

Because the same removal law was used on each mesh, the lower rates could be traced directly to less film in the lower-wall removal region. The useful comparison was therefore how each prepared flow supplied that region. The 2.6M case in Figure 20 provides an initial bulk-flow reference.

Bulk pressure differences varied less across these cases. The calculated area-averaged inlet-to-outlet static pressure difference was 11.0--11.3 kPa, a spread of about 3\% of the mean. This falls within Pointon et al.'s roughly 10--40 kPa range of typical predicted losses. Their example design specified 19 kPa at 11.7 bar and 1875 t/h, compared with the present 11.2 bar and about 711 t/h, so it provides context at different operating conditions (\href{https://publications.mygeoenergynow.org/grc/1028587.pdf}{Pointon et al., 2009, pp.~944--945}).

Using the outlet steam-quality definition in Section 2.2, the bulk vapour and liquid mass-flow rates at the comparison points give an outlet steam quality of 99.32--99.42\%, compared with Pointon's design outlet steam quality of 99.87\%. Appendix A.6 records the calculation. The corresponding liquid fractions, 0.58--0.68\% here and 0.13\% in that design, show how similar steam-quality percentages can hide differences in liquid carryover.

Rizaldy et al.~also found evidence of brine carryover at Wairakei and proposed wall-film entrainment as a contributor, giving the film-transport question practical relevance (\href{https://www.researchgate.net/publication/310818568_LIQUID_CARRYOVER_IN_GEOTHERMAL_STEAM-WATER_SEPARATORS}{Rizaldy et al., 2016, Sections 3--4}).

Similar bulk pressure differences occurred alongside clear changes in film coverage, storage and drainage. Differences in starting fields limit what can be assigned to mesh resolution alone (Appendix A.6). The comparison showed why mesh assessment must also follow the wall-liquid paths: the amount reaching the drain depends on where the film develops and how it moves.

\subsection{Initial design tests}\label{initial-design-tests}

The fourth research question asks how CFD can inform separator design optimisation. The earlier studies identified a specific design problem: liquid captured at the wall could remain high in the vessel instead of reaching the lower removal region. Since providing removal at the base did not resolve that transport problem, the design focus shifted towards guiding wall liquid after capture, rather than treating capture alone as successful separation.

Two features were proposed to address this transport problem: a vertical blocker to interrupt the recirculation of wall liquid before it meets the incoming steam again, and an angled ramp to guide liquid downwards towards the removal region.

The 154k blocker and 213k ramp cases reached 150 ms of EWF development; Appendix A.8 gives the setup details. Figure 21 presents the associated local mixture-velocity vectors at 50 ms.

The blocker interrupts the local vector pattern, while flow beside the ramp follows its inclined edge. These local flow changes agree with the intended roles of the features: interrupting circulation and guiding liquid downwards.

\par\nopagebreak[4]
\noindent\begin{minipage}{\linewidth}
\centering
\begin{minipage}[t]{0.48\linewidth}\centering
\includegraphics[height=40mm,keepaspectratio]{\detokenize{../../experiments/phase-10-proposed-design/figures/manual-vertical-blocker-basic-50ms.png}}
\end{minipage}
\hfill
\begin{minipage}[t]{0.48\linewidth}\centering
\includegraphics[height=40mm,keepaspectratio]{\detokenize{../../experiments/phase-10-proposed-design/figures/manual-angled-ramp-50ms.png}}
\end{minipage}
\captionsetup{hypcap=false}\captionof{figure}[Proposed film blocker and ramp]{Preliminary mixture-velocity vectors around (a) the vertical film blocker and (b) the angled film ramp.}
\label{fig:21}
\end{minipage}\par\medskip

The EWF calculations were too short to show whether these local flow changes increase or reduce separation efficiency. Longer film calculations are needed to assess how each feature changes wall-liquid storage, collection and carryover. These initial tests gave a clear direction for further design work: use a blocker to interrupt circulation or a ramp to guide liquid towards removal.

\subsection{Workflow and automation}\label{workflow-and-automation}

\par\noindent\begin{minipage}{\linewidth}
The preceding investigations required repeated setup changes, long calculations and comparisons between saved states. Automation connected these tasks within a reusable workflow (Figure 22), with OpenAI Codex helping to prepare cases, run Fluent and analyse results. Its main benefits were preserving previous work, carrying out repeated tasks and checking results.

\medskip
\centering
\includegraphics[width=0.98\linewidth,height=0.43\textheight,keepaspectratio]{\detokenize{../../../assets/readme/experiment-workflow.png}}
\captionsetup{hypcap=false}\captionof{figure}[Experiment workflow linking research questions, execution, evidence and subsequent decisions.]{Experiment workflow linking research questions, execution, evidence and subsequent decisions.}
\label{fig:22}
\end{minipage}\par\medskip

The programme included 387 configurations run across both project partners, including failed and partial trials. This programme depended on being able to reuse earlier cases and findings as each investigation developed.



A connected knowledge base and experiment records preserved each test's purpose, assumptions, settings and results. Saved cases and checked procedures gave later tests a starting point, while checkpoints preserved work for recovery after a failure.

Task instructions defined the objectives, allowed changes and stopping conditions. The Model Context Protocol (MCP) connection linked the agent to persistent Fluent sessions through PyFluent, allowing it to inspect settings and issue commands. Scheduled monitoring supported progress checks and recovery during long runs while the host and solver remained available. The workflows and Fluent integration are available in the \href{https://github.com/Shuhei-uni/P4P_knowledgeWiki}{project GitHub repository}.

Checks of Fluent settings, solver logs and saved flow fields allowed the model setup and results to be reviewed. For example, the timestep test in Section 4.8 rejected the 25 µs candidate because of its mass-balance error, despite similar film fields. The mass-balance check tied that decision to numerical evidence.

Scientific judgement remained central to defining the research questions and interpreting the results. Experience gained with Fluent and CFD showed that some investigations added little to understanding the liquid path. Selecting useful tests therefore remained a research decision; automation made it practical to carry them out, preserve the evidence and build on the results.

% Source: draft/04-discussion-conclusions-and-future-work.md
\section{Discussion}\label{discussion}

The study developed the inherited model to follow liquid from entry, through wall capture, to removal. Inlet and mesh changes revealed how liquid gathered near the wall. The absorber and improved startup then made wall-film studies practical. These steps addressed the gaps in Section 1.2 and showed why collection requires both transport towards the drain and removal at the drain.

\subsection{Inlet and mesh representation}\label{inlet-and-mesh-representation}

The split inlet represented the expected concentration of liquid beside the wall as the feed turned into the separator. Complete phase separation at entry remained an idealisation. Compared with mixed feed, it produced a clearer liquid-rich wall band and deeper base accumulation, even though outlet liquid flow changed little (Section 4.1). This showed how representing wall-side entry changed the liquid path inside the vessel.

The move from tetrahedral to poly-hexcore meshes also aimed to improve the representation of swirling flow and liquid near the wall. The spatial comparison showed stronger wall-liquid retention, while the finer meshes showed rising inventory; the 680k continuation eventually failed (Section 4.2). Reduced numerical spreading is a possible explanation for the stronger retention, but differences in inlet and preparation prevent a firm accuracy comparison. The useful finding was that the new cases exposed sustained liquid accumulation that the earlier model had not shown as clearly.

Together, these changes showed where liquid gathered and how storage grew in the closed-bottom domain. The inlet and mesh choices gave a clearer basis for understanding liquid movement and retention within the separator.

\subsection{Liquid removal and numerical stability}\label{liquid-removal-and-numerical-stability}

The retained liquid needed a destination near the base, but tests with the brine pool encountered stability and discharge-control problems. The absorber provided removal directly from the bulk flow, excluding pool motion, pressure response and level control. This separated the question of whether liquid could reach collection from the difficulties of modelling the pool itself.

Liquid removal and reduced-feed startup lowered the reported continuity residual and retained inventory relative to the absorber-free case (Section 4.4). Mesh and loading differed, and common residual scaling was not verified, so the improvement cannot be assigned to the absorber alone. Inventory changed little late in the reduced-feed hold, and the calculation reached full feed. This improved behaviour justified saving a common baseline for wall-treatment tests.

The shortened startup addressed the next practical step: reaching full feed with smaller loading disturbances. With the same continuity normalisation, it reduced the ramp peaks in residual, combined liquid inventory and outlet liquid flow by about 70--76\% (Section 4.7). A separate activation spike occurred before the ramp, and the liquid mass balance remained unresolved. The smoother transition to full feed provided a defined starting field for longer film calculations.

\subsection{Wall--liquid transport and film development}\label{wallliquid-transport-and-film-development}

With removal available at the base, attention turned to liquid remaining higher on the wall. Roughness strongly changed outlet liquid flow and stored mass, but increasing roughness did not always reduce outlet flow. Later work used a commercial-steel reference instead of choosing roughness to minimise carryover. The tests established that wall friction strongly affected the predicted liquid path.

EWF followed what happened after wall capture by representing wall storage and transport, extending the droplet-tracking approach discussed in Section 2.3.1. With the bulk flow held fixed, the film spread and developed sustained drainage through 500 ms. Input still exceeded drainage, so film mass kept growing as outflow began to level out. Tracking both quantities allowed the model to distinguish wall transport from continued storage.

Continued storage prompted release tests, but returning film liquid to droplets did not stop accumulation. A direct lower-wall drain reduced local storage, although some removal replaced existing film outflow. Liquid still accumulated higher on the wall. The drain worked where liquid reached it; guiding liquid down from the upper wall became the design target.

This transport problem also matters in operating separators. Rizaldy et al.~found evidence of brine carryover at Wairakei and proposed release from the wall film as a contributor (Sections 2.4 and 4.10). A lower film mass could therefore mean either more collection or more liquid returning to the steam. The film studies identify what a design comparison must track together: stored liquid, external collection, carryover and pressure loss.

\subsection{Mesh sensitivity of the developed model}\label{mesh-sensitivity-of-the-developed-model}

The developed film model allowed the mesh study to examine whether the coarse-mesh behaviour persisted with refinement. At the same film time, finer meshes showed more distinct film paths, less upper-wall coverage, less stored liquid and lower drainage. The finest completed case also differed from the two intermediate meshes, which had appeared to agree (Section 4.10). The coarse-mesh result therefore did not describe the film response across resolutions.

Because the same removal law was used, lower drain rates followed smaller film inventories in the removal region. They did not by themselves show a less effective drain or less total liquid collected over time. Starting fields differed, and liquid supply and release could also affect the result. Matched preparation and longer histories are needed to separate these effects before claiming mesh independence.

The comparison showed why broad film coverage and strong drainage on one mesh are insufficient to judge collection. Mesh assessment must also follow the liquid paths and the amount reaching the drain.

\subsection{Implications for separator design}\label{implications-for-separator-design}

Effective lower-wall removal alongside continued upper-wall storage gave the design work a clear objective: guide liquid towards the lower drain. A vertical blocker was proposed to interrupt liquid circulating around the vessel, and an angled ramp to guide it downwards. Both features addressed the liquid path identified by the model.

The initial mixture-velocity plots showed local flow changes consistent with these intended roles (Section 4.11). The short calculations did not establish an improvement in collection; that requires longer tests in which bulk flow, film and released droplets interact. Even at this stage, CFD had linked an observed transport problem to specific design changes, giving the next comparisons a physical basis.

\subsection{Overall contribution and remaining limits}\label{overall-contribution-and-remaining-limits}

The main contribution was to make more of the liquid path available for study. The inherited Purnanto-based model was extended to represent removal near the base, film storage and transport, and direct film drainage (Section 2.3.1). These additions allowed the investigation to follow liquid after wall capture and test ways of guiding it towards collection.

Failed and partial trials also shaped the work. Pool-control problems led to simplified removal, while timestep trials showed why short checks did not guarantee stable longer runs. Automation helped repeat setup, checks, recovery and comparison. Scientific judgement guided the choice of research questions and interpretation of the results (Section 4.12).

Computational cost limited film duration and the number of mesh comparisons. Holding the bulk flow fixed made longer film runs practical, but left the changing interaction between bulk flow, film and released droplets unresolved. The combined liquid mass balance, mesh independence and agreement with measurements also remain to be established. These gaps limit quantitative claims about separation efficiency and pressure loss.

The study provided model features and comparison procedures for following liquid delivery, storage, transport and removal together. This connected view identified a specific design problem: guiding upper-wall liquid towards the drain.

\section{Conclusions and future work}\label{conclusions-and-future-work}

\subsection{Conclusions}\label{conclusions}

The study extended the inherited CFD model to investigate liquid transport and collection. The split inlet represented expected wall-side entry, while poly-hexcore meshes aimed to better represent swirling flow and wall liquid. These changes exposed clearer wall accumulation and showed why internal distribution must be assessed alongside outlet flow.

The absorber and revised startup supported wall-treatment studies. The shortened startup reduced the peak continuity residual during the feed ramp by about 70\%. Full convergence remained unresolved, but the smoother loading procedure supplied a practical starting field for longer film calculations.

Roughness changed the liquid path, and EWF made film storage and transport explicit. Drainage developed through 500 ms with the bulk flow held fixed, while film mass kept rising. Direct drainage controlled lower-wall storage, but liquid still accumulated above. Collection therefore requires both transport towards the drain and removal at the drain.

Finer meshes changed film paths, stored mass and drainage. Mesh independence, the combined liquid mass balance and separation performance remain unverified. The comparison identified film development and liquid supply to the drain as key parts of the next assessment.

The blocker and ramp turn this understanding into design ideas. Initial flow patterns supported further testing, although improved collection has not yet been shown.

The work made substantial progress towards the main research goal of developing a numerically consistent and physically credible CFD model of liquid transport and collection. Improved startup reduced loading disturbances, while liquid-removal and wall-film modelling made more of the liquid path available for study. The results linked changes in liquid distribution to changes in storage and drainage, and identified upper-wall retention as a specific design problem. These advances delivered both a more capable model and a deeper understanding of liquid behaviour within the separator, providing a physical basis for designs that guide captured liquid towards collection.

\subsection{Future work}\label{future-work}

The first priority is to run the film for longer and determine whether storage stops growing. Drainage had begun to level out at 500 ms, but film mass was still rising. Longer runs with the bulk flow held fixed should track total and upper-wall film mass, input, release and drainage. A statistically steady film would have no sustained mass growth and balanced mean input and outflow over successive time windows. Timestep and Courant settings should first be checked on a developed film so these longer runs remain practical.

If storage keeps growing, local flow and mass histories should help identify whether high input, weak downward transport or limited outflow sustains it. Liquid released as droplets remains inside the separator. A combined bulk and film mass balance should therefore track inlet flow, external removal and stored mass, with release recorded as an internal transfer.

The same histories should be compared on selected finer meshes, using matched preparation and film settings over longer periods. Their lower film mass and drain rates at about 200 ms may reflect slower development or different liquid supply and release. Following these quantities over time would show whether the responses become similar or remain different, helping separate development time from mesh effects.

The next design tests should compare the blocker and ramp with the unmodified separator under matched feed, film settings and preparation. Bulk flow, film and released droplets must interact so the comparison can track upper-wall storage, external collection and liquid returning to the steam. These tests would establish whether each feature guides liquid towards removal and improves collection.

Finally, pressure loss, outlet liquid carryover and collected liquid flow should be compared with measurements at matched operating conditions. This would test whether the added model features reproduce separator behaviour closely enough to support quantitative design decisions.

\clearpage\ESBfrontsection{Acknowledgements}
I thank Sadiq Zarrouk and Maedeh Amirpour for their supervision and feedback as the modelling questions developed. Discussions with Sadiq also helped shape the wall-liquid design concepts. I thank Andy Tran for his collaboration as my project partner.

\textbf{Contribution statement.} This report presents my simplified-domain inlet, mesh, liquid-removal, startup and wall-film investigations. Andy Tran carried out the separate replication and full-geometry numerical investigations. The total of 387 executed configurations in Section 4.12 covers both partners; it is not a count of my work alone. Supervisor discussions helped shape the research direction and wall-liquid design concepts.

\clearpage\ESBfrontsection{References}\begingroup\setstretch{1.0}\setlength{\parskip}{3pt}\raggedright\interlinepenalty=10000\setlength{\parindent}{-1.27cm}\setlength{\leftskip}{1.27cm}
Ansys, Inc.~(2020). \emph{Ansys Fluent Mosaic technology automatically combines disparate meshes with polyhedral elements for fast, accurate flow resolution} {[}White paper{]}. \url{https://www.ansys.com/content/dam/resource-center/white-paper/ansys-fluent-mosaic-technology-wp.pdf}

Ansys, Inc.~(2025a). Cell zone conditions (§8.2.7, Defining mass, momentum, energy, and other sources). In \emph{Ansys Fluent user's guide} (Release 2025 R2). \url{https://ansyshelp.ansys.com/public/Views/Secured/corp/v252/en/flu_ug/flu_ug_bcs_sec_cell_zones.html\#flu_ug_sec_bc_source}

Ansys, Inc.~(2025b). Film sub-models (§17.2.1). In \emph{Ansys Fluent theory guide} (Release 2025 R2). \url{https://ansyshelp.ansys.com/public/Views/Secured/corp/v252/en/flu_th/flu_th_ewf_sec_film_submodels.html}

Ansys, Inc.~(2025c). Mesh topologies (§7.1, particularly §7.1.3.3, Numerical diffusion). In \emph{Ansys Fluent user's guide} (Release 2025 R2). \url{https://ansyshelp.ansys.com/public/Views/Secured/corp/v252/en/flu_ug/flu_ug_GridTypes.html}

Ansys, Inc.~(2025d). Mixture model theory (§14.4, particularly the overview and relative-velocity formulation). In \emph{Ansys Fluent theory guide} (Release 2025 R2). \url{https://ansyshelp.ansys.com/public/Views/Secured/corp/v252/en/flu_th/flu_th_sec_multiphase_mixture.html}

Ansys, Inc.~(2025e). Modeling Eulerian wall films (Chapter 30, particularly §§30.3--4). In \emph{Ansys Fluent user's guide} (Release 2025 R2). \url{https://ansyshelp.ansys.com/public/views/secured/corp/v252/en/flu_ug/flu_ug_models_wallfilm.html}

Ansys, Inc.~(2025f). Model-specific DEFINE macros (§2.3.45, DEFINE\_SOURCE). In \emph{Ansys Fluent customization manual} (Release 2025 R2). \url{https://ansyshelp.ansys.com/public/Views/Secured/corp/v252/en/flu_udf/flu_udf_ModelSpecificDEFINE.html\#flu_udf_sec_define_source}

Ansys, Inc.~(2025g). Numerical schemes and solution algorithm (§17.4). In \emph{Ansys Fluent theory guide} (Release 2025 R2). \url{https://ansyshelp.ansys.com/public/Views/Secured/corp/v252/en/flu_th/flu_th_ewf_sec_num_schemes.html}

Ansys, Inc.~(2025h). Overview of using the Eulerian wall film model (§30.2). In \emph{Ansys Fluent user's guide} (Release 2025 R2). \url{https://ansyshelp.ansys.com/public/views/secured/corp/v252/en/flu_ug/flu_ug_ewf_sec_overview.html}

Ansys, Inc.~(2025i). Pressure-based solver settings (§37.3, particularly §37.3.1, Choosing the pressure--velocity coupling method). In \emph{Ansys Fluent user's guide} (Release 2025 R2). \url{https://ansyshelp.ansys.com/public/Views/Secured/corp/v252/en/flu_ug/flu_ug_sec_uns_solve_pvel_1.html}

Ansys, Inc.~(2025j). Solution strategies for multiphase modeling (§27.8, particularly §§27.8.1.2--1.3 and 27.8.1.7). In \emph{Ansys Fluent user's guide} (Release 2025 R2). \url{https://ansyshelp.ansys.com/public/Views/Secured/corp/v252/en/flu_ug/flu_ug_sec_multiphase_solution.html}

Ansys, Inc.~(2025k). Wall roughness effects in turbulent wall-bounded flows (§4.18.5). In \emph{Ansys Fluent theory guide} (Release 2025 R2). \url{https://ansyshelp.ansys.com/public/Views/Secured/corp/v252/en/flu_th/x1-5720008.164.html}

Ansys, Inc.~(2025l). Swirling and rotating flows (§10.3, particularly §§10.3.1.2, 10.3.2--10.3.3 and 10.3.4.2.2). In \emph{Ansys Fluent user's guide} (Release 2025 R2). \url{https://ansyshelp.ansys.com/public/Views/Secured/corp/v252/en/flu_ug/flu_ug_uns_sec_models_swirl.html}

Ansys, Inc.~(2025m). Convergence and stability (§37.22.1, Judging convergence). In \emph{Ansys Fluent user's guide} (Release 2025 R2). \url{https://ansyshelp.ansys.com/public/views/secured/corp/v252/en/flu_ug/flu_ug_sec_solve_conv_stab.html}

El Haj Assad, M., Bani-Hani, E., \& Khalil, M. (2017). Performance of geothermal power plants (single, dual, and binary) to compensate for LHC-CERN power consumption: Comparative study. \emph{Geothermal Energy, 5}, Article 17. \url{https://doi.org/10.1186/s40517-017-0074-z}

International Association for the Properties of Water and Steam. (2012). \emph{Revised release on the IAPWS industrial formulation 1997 for the thermodynamic properties of water and steam} (R7-97{[}2012{]}). \url{https://iapws.org/technical-guidance/release/IF97-Rev}

National Aeronautics and Space Administration. (2021a, February 10). \emph{Examining iterative convergence}. NPARC Alliance CFD verification and validation tutorial. \url{https://www.grc.nasa.gov/www/wind/valid/tutorial/iterconv.html}

National Aeronautics and Space Administration. (2021b, February 10). \emph{Validation assessment}. NPARC Alliance CFD verification and validation tutorial. \url{https://www.grc.nasa.gov/www/wind/valid/tutorial/valassess.html}

National Aeronautics and Space Administration. (2021c, February 10). \emph{Verification assessment}. NPARC Alliance CFD verification and validation tutorial. \url{https://www.grc.nasa.gov/www/wind/valid/tutorial/verassess.html}

Pointon, A. R., Mills, T. D., Seil, G. J., \& Zhang, Q. (2009). Computational fluid dynamic techniques for validating geothermal separator sizing. \emph{GRC Transactions, 33}, 943--948. \url{https://publications.mygeoenergynow.org/grc/1028587.pdf}

Purnanto, M. H., Zarrouk, S. J., \& Cater, J. E. (2013). CFD modelling of two-phase flow inside geothermal steam--water separators. \emph{IPENZ Transactions, 40}, 1--10. \url{https://www.researchgate.net/publication/269519626_CFD_MODELLING_OF_TWO-PHASE_FLOW_INSIDE_GEOTHERMAL_STEAM-WATER_SEPARATORS}

Rizaldy, Zarrouk, S. J., \& Morris, C. (2016). Liquid carryover in geothermal steam--water separators. In \emph{Proceedings of the 38th New Zealand Geothermal Workshop}, Auckland, New Zealand. \url{https://www.researchgate.net/publication/310818568_LIQUID_CARRYOVER_IN_GEOTHERMAL_STEAM-WATER_SEPARATORS}

Skoog, E. (2020). \emph{CFD annular flow modelling based on a three-field approach} {[}Master's thesis, Luleå University of Technology{]}. DiVA. \url{https://www.diva-portal.org/smash/get/diva2:1452100/FULLTEXT02.pdf}

U.S. Energy Information Administration. (n.d.). \emph{Geothermal power plants}. \url{https://www.eia.gov/energyexplained/geothermal/geothermal-power-plants.php}

Zarrouk, S. J., \& Purnanto, M. H. (2015). Geothermal steam--water separators: Design overview. \emph{Geothermics, 53}, 236--254. \url{https://doi.org/10.1016/j.geothermics.2014.05.009}

\endgroup
\clearpage\ESBappendix
\renewcommand{\thefigure}{\thesection\arabic{figure}}
\renewcommand{\thetable}{\thesection\arabic{table}}
\renewcommand{\theHtable}{\thetable}
\begin{landscape}
\section{Supporting calculations and numerical checks}\label{appendix-a.-supporting-calculations-and-numerical-checks}

\begingroup\footnotesize\setstretch{1.0}
This appendix gives the model settings, calculations and numerical checks that support Sections 3--5. It explains how the comparisons were made and how the reported quantities were calculated.

\subsection{Model controls and comparison basis}\label{a.1.-model-controls-and-comparison-basis}

Table A1a identifies each comparison and its limits; Table A1b gives the local controls. Table A1c records configuration-specific film and particle inputs.
\endgroup

\begingroup\setstretch{1.0}\fontsize{9}{10.5}\selectfont\setlength{\tabcolsep}{3pt}\setlength{\LTpre}{3pt}
{
\renewcommand{\thetable}{A1a}
\begin{longtable}[]{@{}
  >{\raggedright\arraybackslash}p{(\linewidth - 12\tabcolsep) * \real{0.1300}}
  >{\raggedright\arraybackslash}p{(\linewidth - 12\tabcolsep) * \real{0.0900}}
  >{\raggedright\arraybackslash}p{(\linewidth - 12\tabcolsep) * \real{0.1300}}
  >{\raggedright\arraybackslash}p{(\linewidth - 12\tabcolsep) * \real{0.1600}}
  >{\raggedright\arraybackslash}p{(\linewidth - 12\tabcolsep) * \real{0.1700}}
  >{\raggedright\arraybackslash}p{(\linewidth - 12\tabcolsep) * \real{0.1500}}
  >{\raggedright\arraybackslash}p{(\linewidth - 12\tabcolsep) * \real{0.1700}}@{}}
\caption{Meshes and comparison basis for the reported investigations.}\label{tab:A1a}\tabularnewline
\toprule\noalign{}
\begin{minipage}[b]{\linewidth}\raggedright
Investigation
\end{minipage} & \begin{minipage}[b]{\linewidth}\raggedright
Mesh
\end{minipage} & \begin{minipage}[b]{\linewidth}\raggedright
Pressure--velocity scheme
\end{minipage} & \begin{minipage}[b]{\linewidth}\raggedright
Initialisation / starting field
\end{minipage} & \begin{minipage}[b]{\linewidth}\raggedright
Run length / comparison window
\end{minipage} & \begin{minipage}[b]{\linewidth}\raggedright
Main change
\end{minipage} & \begin{minipage}[b]{\linewidth}\raggedright
Other factors affecting interpretation
\end{minipage} \\
\midrule\noalign{}
\endfirsthead
\toprule\noalign{}
\begin{minipage}[b]{\linewidth}\raggedright
Investigation
\end{minipage} & \begin{minipage}[b]{\linewidth}\raggedright
Mesh
\end{minipage} & \begin{minipage}[b]{\linewidth}\raggedright
Pressure--velocity scheme
\end{minipage} & \begin{minipage}[b]{\linewidth}\raggedright
Initialisation / starting field
\end{minipage} & \begin{minipage}[b]{\linewidth}\raggedright
Run length / comparison window
\end{minipage} & \begin{minipage}[b]{\linewidth}\raggedright
Main change
\end{minipage} & \begin{minipage}[b]{\linewidth}\raggedright
Other factors affecting interpretation
\end{minipage} \\
\midrule\noalign{}
\endhead
\bottomrule\noalign{}
\endlastfoot
Inlet arrangement (4.1) & 60k & Coupled & Fresh Hybrid Initialisation for each case & 10,000 updates; final 500 & Mixed / split inlet; absorber off & Idealised inlet distributions, not measured profiles \\
Accumulation (4.2) & 342k, 680k, 997k; inherited tetrahedral mesh & SIMPLE in finer-mesh tests & Fresh Hybrid Initialisation after mesh replacement & Comparisons up to 3000 iterations; 680k run continued to failure & Absorber-free finer meshes & Tetrahedral comparison also changes inlet, preparation and controls \\
Liquid removal and startup (4.3--4.4) & 60k & SIMPLE, then Coupled & Hybrid; reduced-feed hold, ramp and saved full-feed field & Reduced-feed endpoint N1580; wall-test baseline N5586 & Absorber and loading procedure & 342k reference differs in mesh, feed and residual scaling \\
Wall treatment (4.6) & 60k & Coupled & Independent restarts from N5586 & 3000 updates; final-500 means & Roughness with EWF off; EWF with smooth walls & Developing baseline; several film settings change together \\
Early activation / film development (4.7--4.8) & 60k & Coupled; bulk flow later held fixed & Retain N1580 bulk flow; initialise dry film & 500 hold + 2000 ramp + 1000 full-feed updates; film to 500 ms & Earlier combined activation; then film-only development & Several startup settings change together; bulk flow held fixed during film development \\
Release and drainage (4.9) & 60k & Coupled preparation; bulk flow held fixed for drain pair & Retain developed fields; identical lower-film extension in OFF/ON pair & Matched 15 ms drain test; separate longer continuation & Release mechanisms; direct film drain & Release study differs from 500 ms case; only drain pair isolates direct removal \\
Mesh sensitivity (4.10) & 60k, 342k, 680k, 997k & Coupled preparation; bulk flow held fixed & Mapped fields and mesh-specific startup & Common film time of 195 ms & Resolution with steel roughness and direct film drain & Different bulk preparation and starting film times \\
Proposed designs (4.11) & 154k / 213k & Coupled preparation; bulk flow held fixed & Separate prepared design fields & 150 ms total; final 100 ms at 10 µs & Blocker / angled ramp; basic EWF & Geometry, mesh and starting fields differ \\
\end{longtable}
}
\endgroup
\end{landscape}

The selected comparisons used individual runs, so within-window fluctuations do not measure run-to-run uncertainty. Fixed run lengths and common film times define comparison points, not steady state. Exploratory turbulence trials informed model selection; a matched turbulence-model comparison was not completed for the film findings.

The 60k label denotes the 60,964-cell poly-hexcore development mesh; N denotes solver iteration. The inlet comparison kept the same mesh regions with the absorber inactive; removal and wall-treatment studies activated liquid removal near the base. Mixed and split inlets shared the Table A1a setup, while wall-treatment tests reused the flow field at iteration 5586. Finer absorber-free tests started from Hybrid Initialisation after mesh replacement. The film comparisons in Appendix A.6 instead used prepared bulk flows held fixed during film development.

\begingroup\small\setlength{\tabcolsep}{3pt}
{
\renewcommand{\thetable}{A1b}
\begin{longtable}[]{@{}
  >{\raggedright\arraybackslash}p{(\linewidth - 4\tabcolsep) * \real{0.2200}}
  >{\raggedright\arraybackslash}p{(\linewidth - 4\tabcolsep) * \real{0.4900}}
  >{\raggedright\arraybackslash}p{(\linewidth - 4\tabcolsep) * \real{0.2900}}@{}}
\caption{Controls needed to interpret the selected comparisons.}\label{tab:A1b}\tabularnewline
\toprule\noalign{}
\begin{minipage}[b]{\linewidth}\raggedright
Comparison
\end{minipage} & \begin{minipage}[b]{\linewidth}\raggedright
Controls and reporting basis
\end{minipage} & \begin{minipage}[b]{\linewidth}\raggedright
Interpretation limit
\end{minipage} \\
\midrule\noalign{}
\endfirsthead
\toprule\noalign{}
\begin{minipage}[b]{\linewidth}\raggedright
Comparison
\end{minipage} & \begin{minipage}[b]{\linewidth}\raggedright
Controls and reporting basis
\end{minipage} & \begin{minipage}[b]{\linewidth}\raggedright
Interpretation limit
\end{minipage} \\
\midrule\noalign{}
\endhead
\bottomrule\noalign{}
\endlastfoot
Bulk model (Section 3.2) & Mixture; RNG \(k\)--\(\varepsilon\), swirl treatment and standard wall functions; Manninen slip, Schiller--Naumann drag, 10 µm liquid diameter; gravity \((0,-9.81,0)\) m/s²; energy off & No flashing or condensation; applicability to concentrated wall liquid unverified \\
Full-feed control & 60,964 cells; liquid/vapour feed 116.92/80.69 kg/s; outlet 1.12 MPa absolute; Coupled/Global Time Step; PRESTO! pressure, QUICK volume fraction, second-order momentum and \(\varepsilon\), first-order \(k\) & Control window: iterations 4586--5586; a fixed horizon does not establish convergence \\
Inlet and mesh comparisons (Sections 4.1--4.2) & Absorber off. Mixed/split Coupled cases share coarse mesh and total phase feed; fresh Hybrid Initialisation; 10,000 iterations; reporting window 9500--10,000 & SIMPLE tests also change pseudo-time and \(k\) scheme; tetrahedral comparison changes mesh and inlet placement \\
Roughness and early film tests (Section 4.6) & Independent restarts from iteration 5586; 3000 updates, final-500 means. Roughness: EWF off. Film: smooth walls, accretion and coupled film equations on, 10 µs step, momentum feedback off & Starting flow still developing; several film settings change together \\
Combined continuation & Original film fields at iteration 13586; contact-based absorber, 0.5 mm roughness at \(C_s=0.5\); 1 µs film step, then adaptive continuation to iteration 45606; final accretion/drainage 82.152/73.956 kg/s (9.977\% deficit) & Several settings change; this is a development sequence, not an isolated treatment effect \\
Shortened startup (Section 4.7) & Reuse iteration-1580 fields; activate the contact-based absorber, Coupled scheme, 0.5 mm roughness and dry EWF; 500 updates at 25\% feed, 2000-update ramp, 1000 full-feed updates; 1 µs film step & Compare ramp progress separately from the activation hold \\
Film development (Section 4.8) & Hold bulk flow fixed after iteration 5080; alternative implicit film algorithm; accepted film time to 500 ms & Film develops under a fixed bulk flow; no evolving whole-separator mass balance \\
Release and drainage (Section 4.9) & Steel reference: 0.045 mm roughness, \(C_s=0.5\), smooth bottom. Stripping and edge separation enabled; film-to-bulk momentum feedback off. Drain OFF/ON cases use identical starting fields and the same film domain extended onto the lower wall & Different preparation from the 500 ms film test; only the matched drain pair isolates direct removal \\
\end{longtable}
}
\endgroup

The full-feed control used 2.11\% inlet turbulence intensity and liquid/vapour hydraulic diameters of 0.01338/0.72061 m. Outlet backflow contained no liquid and used 2.11\% turbulence intensity and a hydraulic diameter of 0.875936 m. Pressure and momentum relaxation factors were 0.5; explicit volume-fraction relaxation was reduced from 0.5 in the control to 0.1 in the contact-based absorber and shortened-startup cases.

\subsubsection{Film and particle inputs by configuration}\label{film-and-particle-inputs-by-configuration}

\begingroup\small\setlength{\tabcolsep}{3pt}
{
\renewcommand{\thetable}{A1c}
\begin{longtable}[]{@{}
  >{\raggedright\arraybackslash}p{(\linewidth - 2\tabcolsep) * \real{0.2300}}
  >{\raggedright\arraybackslash}p{(\linewidth - 2\tabcolsep) * \real{0.7700}}@{}}
\caption{Film and particle settings for each test group.}\label{tab:A1c}\tabularnewline
\toprule\noalign{}
\begin{minipage}[b]{\linewidth}\raggedright
Input
\end{minipage} & \begin{minipage}[b]{\linewidth}\raggedright
Setting and applicable tests
\end{minipage} \\
\midrule\noalign{}
\endfirsthead
\toprule\noalign{}
\begin{minipage}[b]{\linewidth}\raggedright
Input
\end{minipage} & \begin{minipage}[b]{\linewidth}\raggedright
Setting and applicable tests
\end{minipage} \\
\midrule\noalign{}
\endhead
\bottomrule\noalign{}
\endlastfoot
Early film tests and shortened startup (4.6--4.8) & Phase accretion, film momentum and coupled film mass/momentum solution enabled; film-to-bulk momentum feedback disabled. Initial film step 10 µs in the test and 1 µs in shortened startup. The 500 ms continuation used the alternative implicit algorithm; local timestep checks are in A.4. \\
Release tests (4.9) & Film pressure gradient, spreading, surface tension, gravity, surface shear, phase accretion, DPM collection, splash, stripping and edge separation enabled. Film momentum and coupled film solution enabled; film-to-bulk momentum feedback disabled. Film surface tension 0.07194 N/m; energy and phase change not modelled. \\
Specified mesh-comparison settings (4.10) & Bulk and film surface tension 0.04041 N/m; pressure gradient, spreading, gravity, shear, momentum advection, accretion, DPM collection/splash, stripping and edge separation enabled. Film-to-bulk momentum feedback disabled. Sequential film mass/momentum solution, original first-order implicit time scheme, first-order spatial schemes, up to 100 inner iterations and residual tolerance \(10^{-5}\). Curvature smoothing: level 2, factor 0.5. Identical use of these settings was not confirmed for every mesh. \\
Film-edge treatment & Early film and basic-design cases allowed film to leave at the edge of the upper-wall film domain. Direct-drain and mesh tests extended film onto the lower collector wall and removed mass and momentum there with a 1.5 ms depletion time. Film-edge outflow and direct removal were counted separately (A.5). \\
Contact angle in the direct-drain test & Contact-angle force was disabled on the lower film wall; an upper-wall contact-angle value was not established. The separation angle below controls film release and is distinct from the wetting contact angle. \\
Release coefficients in the direct-drain test & Stripping diameter constant 0.14, mass constant 0.5, stripping constant 100 and beta 0.3. Separation angle 0.349066 rad (20°), separation constant 0 and damping 1; random edge separation disabled. These coefficients are specific to this configuration. \\
Basic design configuration (4.11) & Upper-wall EWF; lower wall outside film domain; film-edge outflow; direct film drain disabled. Roughness 0.5 mm, \(C_s=0.5\), smooth physical bottom. Stripping and edge separation disabled; 10 µs step and bulk flow held fixed for the 50--150 ms continuations. \\
\end{longtable}
}
\endgroup

The direct-drain configuration included six diagnostic particle injections at the steam inlet, with diameters of 5.63, 28.14, 56.27, 112.54, 168.811 and 348.88 µm. Each supplied only \(10^{-20}\) kg/s at velocity \((27.118,0,0)\) m/s, with spherical drag, no turbulent dispersion and fixed injection positions. This negligible feed was used for numerical checks, not to represent the inlet liquid distribution. Droplets released by stripping or edge separation were a separate source. The mesh-comparison procedure specified particle-tracking updates every 20 µs of film time, with the number of film steps between updates adjusted accordingly. These particle diameters are separate from the 10 µm diameter used in the bulk slip model (Section 3.2).

\subsection{Feed composition and inlet-strip calculation}\label{a.2.-feed-composition-and-inlet-strip-calculation}

\par\noindent\begin{minipage}{\linewidth}
The phase mass-flow rates in Section 3.3 follow from the feed enthalpy. Purnanto et al.~(2013), Table 1 gives \(h=1600\) kJ/kg, \(\dot m=197.61\) kg/s and 11.2 bar absolute, with saturated liquid/vapour enthalpies of 784.66/2781.46 kJ/kg:

\[
x_V=\frac{1600-784.66}{2781.46-784.66}=0.408323,
\qquad \dot m_V=80.6888,\quad \dot m_L=116.9212\ \mathrm{kg/s}.
\tag{A1}
\]

The inlet boundary conditions used the rounded values 80.69 and 116.92 kg/s. Enthalpy set the feed composition; phase change was not solved within the domain.
\end{minipage}\par\medskip

The strip width in Section 4.1 follows from the phase volumetric-flow rates, \(Q_i=\dot m_i/\rho_i\). The original split-inlet design assumed a common velocity \(U=(Q_L+Q_V)/A\) and assigned each phase an area \(A_i=Q_i/U\). With \(A=0.724^2=0.524176\) m², the liquid-strip width was \(w_L=A_L/0.724\).

\begin{table}[!htbp]
\centering
\renewcommand{\thetable}{A2}
\caption{Inlet sizing from the original design densities and the densities used in the model.}\label{tab:A2}
\begingroup\small\setlength{\tabcolsep}{3pt}
{
\begin{tabular}{@{}
  >{\raggedright\arraybackslash}p{(\linewidth - 4\tabcolsep) * \real{0.4500}}
  >{\raggedleft\arraybackslash}p{(\linewidth - 4\tabcolsep) * \real{0.2750}}
  >{\raggedleft\arraybackslash}p{(\linewidth - 4\tabcolsep) * \real{0.2750}}@{}}
\toprule\noalign{}
\begin{minipage}[b]{\linewidth}\raggedright
Quantity
\end{minipage} & \begin{minipage}[b]{\linewidth}\raggedleft
Original inlet sizing
\end{minipage} & \begin{minipage}[b]{\linewidth}\raggedleft
Model boundary conditions
\end{minipage} \\
\midrule\noalign{}
Liquid / vapour density (kg/m³) & 881.77 / 5.73 & 881.211 / 5.79743 \\
Liquid / vapour volumetric flow (m³/s) & 0.132597 / 14.0820 & 0.132681 / 13.9182 \\
Liquid inlet area (m²) & 0.00488962 & 0.00488992, supplied mesh \\
Liquid-strip width (mm) & 6.75362 & 6.75403 \\
Combined volumetric-flow speed (m/s) & 27.1180 & 26.8057 \\
Liquid / vapour normal speed (m/s) & 27.1180 / 27.1180 & 27.1336 / 26.8026 \\
\bottomrule\noalign{}
\end{tabular}
}
\endgroup
\end{table}

The model retained the supplied inlet areas and phase mass-flow rates. The nominal 26.81 m/s operating point therefore did not impose exactly equal phase velocities. The inlet-comparison tests instead scaled both phase mass-flow rates together: their 26.81 m/s reference used 116.9387/80.7029 kg/s liquid/vapour, slightly different from the original boundary values.

\subsection{Absorber laws and liquid mass balance}\label{a.3.-absorber-laws-and-liquid-mass-balance}

\par\noindent\begin{minipage}{\linewidth}
The contact-based absorber in Section 4.3 removes liquid from 715 cells whose centroids satisfy \(y\leq0.10\) m:

\[
S_L=-\rho_L\max(\alpha_L,0)/\tau,\qquad \tau=10^{-5}\ \mathrm{s}.
\tag{A2}
\]

It removes liquid momentum using liquid velocity, includes associated \(k\) and \(\varepsilon\) sources, and has no direct vapour mass sink. With \(M_c\) the absorber's liquid mass, calculated removal \(M_c/\tau\) matched the mass sink applied by Fluent throughout all 3500 shortened-startup iterations. The source-law check in Section 4.7 therefore confirmed that the absorber removed liquid at the specified rate. The calculated and applied rates represent the same removal and must be counted only once in a mass balance.
\end{minipage}\par\medskip

The depletion time is assumed, not a measured residence time. The 100-iteration strength tests shared the starting field, liquid-velocity momentum removal and 1 µs film step.

\begin{table}[!htbp]
\centering
\renewcommand{\thetable}{A3a}
\caption{Effect of absorber depletion time in the short strength tests.}\label{tab:A3a}
\begingroup\small\setlength{\tabcolsep}{3pt}
{
\begin{tabular}{@{}
  >{\raggedleft\arraybackslash}p{(\linewidth - 6\tabcolsep) * \real{0.2500}}
  >{\raggedleft\arraybackslash}p{(\linewidth - 6\tabcolsep) * \real{0.2500}}
  >{\raggedleft\arraybackslash}p{(\linewidth - 6\tabcolsep) * \real{0.2500}}
  >{\raggedleft\arraybackslash}p{(\linewidth - 6\tabcolsep) * \real{0.2500}}@{}}
\toprule\noalign{}
\begin{minipage}[b]{\linewidth}\raggedleft
Depletion time (µs)
\end{minipage} & \begin{minipage}[b]{\linewidth}\raggedleft
Final absorber liquid mass (g)
\end{minipage} & \begin{minipage}[b]{\linewidth}\raggedleft
Maximum liquid volume fraction (\%)
\end{minipage} & \begin{minipage}[b]{\linewidth}\raggedleft
Final removal rate (kg/s)
\end{minipage} \\
\midrule\noalign{}
100 & 2.780 & 0.05010 & 27.80 \\
10 & 0.496 & 0.01978 & 49.59 \\
1 & 0.159 & 0.01302 & 158.81 \\
\bottomrule\noalign{}
\end{tabular}
}
\endgroup
\end{table}

The short tests did not establish an optimum or steady removal rate. At 1 µs, the small further decrease in maximum liquid fraction came with greater source variability and numerical sensitivity. The 10 µs depletion time was retained.

\par\noindent\begin{minipage}{\linewidth}
The startup in Section 4.4 and the first wall-treatment tests used a throughput-controlled absorber. Its total removal targeted the inlet liquid mass-flow rate and was distributed by liquid content within the absorber volume \(V_c\):

\[
S_{L,\mathrm{throughput}}=-\frac{\dot m_{L,\mathrm{in}}\alpha_L}
{\max(\int_{V_c}\alpha_L\,\mathrm dV,10^{-6}\ \mathrm{m^3})}.
\tag{A3}
\]
\end{minipage}\par\medskip

\begin{table}[!htbp]
\centering
\renewcommand{\thetable}{A3b}
\caption{Full-feed-control mass balance at iteration 5586, EWF off.}\label{tab:A3b}
\begingroup\small\setlength{\tabcolsep}{3pt}
{
\begin{tabular}{@{}
  >{\raggedright\arraybackslash}p{(\linewidth - 4\tabcolsep) * \real{0.3333}}
  >{\raggedleft\arraybackslash}p{(\linewidth - 4\tabcolsep) * \real{0.3333}}
  >{\raggedleft\arraybackslash}p{(\linewidth - 4\tabcolsep) * \real{0.3333}}@{}}
\toprule\noalign{}
\begin{minipage}[b]{\linewidth}\raggedright
Account
\end{minipage} & \begin{minipage}[b]{\linewidth}\raggedleft
Boundary plus applied source (kg/s)
\end{minipage} & \begin{minipage}[b]{\linewidth}\raggedleft
Remainder / relevant feed (\%)
\end{minipage} \\
\midrule\noalign{}
Liquid & −24.3344 & −20.8129 \\
Vapour & +0.4391 & +0.5441 \\
Mixture (solver report) & −23.8923 & −12.0906 \\
\bottomrule\noalign{}
\end{tabular}
}
\endgroup
\end{table}

Boundary flux is positive inwards and the absorber source is negative. Denominators are 116.92, 80.69 and 197.61 kg/s, respectively. These steady-solver equation remainders are not measured physical storage rates. Their magnitude limits the wall studies to model responses rather than verified separation performance.

During shortened-startup iterations 4581--5080, inlet minus outlet minus applied contact removal was \(116.9200-3.3012-63.0140=50.6048\) kg/s, before EWF transfer. Film accretion was 80.7365 kg/s and storage was 80.7364 kg/s over the same 0.5 ms of accepted film time. Almost all film input was stored. These accounts require compatible bulk--film transfer terms before they can establish whole-separator closure.

For the reduced-feed comparison in Section 4.4, \(0.20837/0.00032056=650.02\) gives the continuity-residual ratio. It compares the minimum during the reduced-feed hold with the absorber-free case's endpoint; common residual scaling was not verified. The inventory reduction is \(100(1-31.357136/453.612255)=93.09\%\), comparing the two cases rather than depletion within one run. Within the reduced-feed hold itself, the estimated inventory fall was only 0.30\% over iterations 1080--1580, indicating little further change in stored liquid.

For the startup-ramp comparison in Section 4.7, both cases used the same continuity normalisation, 386.045961. The comparison covers iterations 1581--3580 of the original startup and 2081--4080 of the shortened startup. This common scaling supports the 70.15\% reduction in peak continuity residual during the ramp.

\subsection{Film timestep and solver limits}\label{a.4.-film-timestep-and-solver-limits}

The timestep check in Section 4.8 tested whether a larger step preserved both the film fields and film mass balance. Identical fields at iteration 13390 and 75.662 ms were advanced by 2.5 ms with bulk flow held fixed and the alternative implicit algorithm. Field differences compare the same wall facets against a 2.5 µs reference. Film mass-balance error is \(100|\Delta M_f+\Delta D_f-I_A|/|I_A|\), where \(\Delta M_f\) is the change in stored film mass, \(\Delta D_f\) is the drained mass and \(I_A\) is accretion integrated over accepted film time.

\begin{table}[!htbp]
\centering
\renewcommand{\thetable}{A4}
\caption{Film-field differences and mass-balance errors over the local 2.5 ms check.}\label{tab:A4}
\begingroup\small\setlength{\tabcolsep}{3pt}
{
\begin{tabular}{@{}
  >{\raggedright\arraybackslash}p{(\linewidth - 6\tabcolsep) * \real{0.5200}}
  >{\raggedleft\arraybackslash}p{(\linewidth - 6\tabcolsep) * \real{0.1600}}
  >{\raggedleft\arraybackslash}p{(\linewidth - 6\tabcolsep) * \real{0.1600}}
  >{\raggedleft\arraybackslash}p{(\linewidth - 6\tabcolsep) * \real{0.1600}}@{}}
\toprule\noalign{}
\begin{minipage}[b]{\linewidth}\raggedright
Check
\end{minipage} & \begin{minipage}[b]{\linewidth}\raggedleft
Local limit (\%)
\end{minipage} & \begin{minipage}[b]{\linewidth}\raggedleft
20 µs
\end{minipage} & \begin{minipage}[b]{\linewidth}\raggedleft
25 µs
\end{minipage} \\
\midrule\noalign{}
Summed absolute facet-mass difference / reference film mass & 1 & 0.00285 & 0.00757 \\
Reference-mass-weighted velocity difference / reference velocity norm & 2 & 0.00207 & 0.00343 \\
Maximum-thickness difference / reference maximum & 2 & 0.000115 & 0.000150 \\
Drainage-rate difference / reference accretion & 1 & 0.00283 & 0.00364 \\
Candidate film mass-balance error & 0.1 & 0.04743 & 0.13391 \\
Reference film mass-balance error & 0.1 & 0.00958 & 0.00958 \\
Decision & All checks required & Pass & Fail: mass balance \\
\bottomrule\noalign{}
\end{tabular}
}
\endgroup
\end{table}

Both timesteps produced finite film fields. The facet-mass difference was normalised by the existing reference film mass, which grew from 3.8635 to 3.9176 kg; it is not the relative error in newly accumulated liquid. The separate mass-balance check tests whether stored mass changes by input minus drainage.

The original film solver required final thickness and both momentum residuals \(\leq10^{-5}\). All 2000 shortened-startup ramp updates met this tolerance, but 15 of the final 100 full-feed-hold updates failed. The alternative implicit solver supplied no inner residuals in these tests. A small film mass-balance error therefore cannot establish inner convergence or a stationary film. A stationary film and a converged whole-separator solution remain unverified.

The check covers 2.5 ms, not timestep independence at 500 ms. Both candidates reproduced similar fields; the mass-balance limit determined why 20 µs passed and 25 µs failed.

\subsection{Film storage and net external removal}\label{a.5.-film-storage-and-net-external-removal}

Table A5 explains the continued film growth in Section 4.8: drainage remained below accretion.

\begin{table}[!htbp]
\centering
\renewcommand{\thetable}{A5}
\caption{Film mass and late-window rates supporting Section 4.8.}\label{tab:A5}
\begingroup\small\setlength{\tabcolsep}{3pt}
{
\begin{tabular}{@{}
  >{\raggedleft\arraybackslash}p{(\linewidth - 10\tabcolsep) * \real{0.1300}}
  >{\raggedleft\arraybackslash}p{(\linewidth - 10\tabcolsep) * \real{0.1500}}
  >{\raggedleft\arraybackslash}p{(\linewidth - 10\tabcolsep) * \real{0.1800}}
  >{\raggedleft\arraybackslash}p{(\linewidth - 10\tabcolsep) * \real{0.1800}}
  >{\raggedleft\arraybackslash}p{(\linewidth - 10\tabcolsep) * \real{0.1800}}
  >{\raggedleft\arraybackslash}p{(\linewidth - 10\tabcolsep) * \real{0.1800}}@{}}
\toprule\noalign{}
\begin{minipage}[b]{\linewidth}\raggedleft
Film time (ms)
\end{minipage} & \begin{minipage}[b]{\linewidth}\raggedleft
Film mass (kg)
\end{minipage} & \begin{minipage}[b]{\linewidth}\raggedleft
Accretion (kg/s)
\end{minipage} & \begin{minipage}[b]{\linewidth}\raggedleft
Drainage (kg/s)
\end{minipage} & \begin{minipage}[b]{\linewidth}\raggedleft
Storage (kg/s)
\end{minipage} & \begin{minipage}[b]{\linewidth}\raggedleft
Drainage deficit (\%)
\end{minipage} \\
\midrule\noalign{}
265.184 & 6.6883 & 81.121 & 68.177 & 12.954 & 15.96 \\
500.000 & 9.5444 & 81.121 & 69.600 & 11.528 & 14.20 \\
\bottomrule\noalign{}
\end{tabular}
}
\endgroup
\end{table}

The first row uses the latest complete 1000-update window; the second uses 490--500 ms, with partial steps weighted by interval overlap. Storage is \(\Delta M_f/\Delta t_f\); the drainage deficit is \(100(\overline A-\overline D)/|\overline A|\), where \(\overline A\) and \(\overline D\) are mean accretion and drainage rates. Drainage changed slowly, but positive storage shows that the film was still growing.

At 500 ms, the unrounded rates give input minus drainage of 11.52160 kg/s, compared with storage of 11.52797 kg/s. The difference is 0.00637 kg/s, or 0.00785\% of input, over the final reporting window. This checks the film mass balance only; the calculated accumulation closely matches the excess input over drainage.

The release test in Section 4.9 added stripping and edge separation, yet film mass still rose from 0.263809 to 5.907035 kg over 50 ms with bulk flow held fixed. Release alone therefore did not stop accumulation in that test.

\par\noindent\begin{minipage}{\linewidth}
The direct film drain addressed liquid that reached the lower wall. It used a separate mass sink \(s_m=-K_d\rho_L\max(h_f,0)\) and momentum source \(s_m\boldsymbol u_f\), where \(h_f\) and \(\boldsymbol u_f\) are film thickness and velocity, with

\[
K_d=\min(1/\tau_f,0.01/\Delta t_{\mathrm{refresh}})=666.667\ \mathrm{s^{-1}},
\quad \tau_f=1.5\ \mathrm{ms},\quad\Delta t_{\mathrm{refresh}}=15\ \mathrm{\mu s}.
\tag{A4}
\]

An isolated drain check with bulk flow held fixed removed 31.9363 g from an initially imposed 50.3752 g film over 1.5 ms; Drain OFF changed by only 0.0000783 g. Fluent's film-edge outflow excludes the direct film drain, so direct removal was integrated separately.
\end{minipage}\par\medskip

\begin{table}[!htbp]
\centering
\renewcommand{\thetable}{A6}
\caption{Matched 15 ms direct-drain comparison, initially 9.7526 kg total film.}\label{tab:A6}
\begingroup\small\setlength{\tabcolsep}{3pt}
{
\begin{tabular}{@{}
  >{\raggedright\arraybackslash}p{(\linewidth - 4\tabcolsep) * \real{0.5600}}
  >{\raggedleft\arraybackslash}p{(\linewidth - 4\tabcolsep) * \real{0.2200}}
  >{\raggedleft\arraybackslash}p{(\linewidth - 4\tabcolsep) * \real{0.2200}}@{}}
\toprule\noalign{}
\begin{minipage}[b]{\linewidth}\raggedright
Quantity
\end{minipage} & \begin{minipage}[b]{\linewidth}\raggedleft
Drain OFF
\end{minipage} & \begin{minipage}[b]{\linewidth}\raggedleft
Drain ON
\end{minipage} \\
\midrule\noalign{}
Final lower-wall film (kg) & 0.18752 & 0.05565 \\
Final total film (kg) & 11.1726 & 10.9829 \\
Integrated direct removal (kg) & 0 & 0.5088 \\
Integrated film-edge outflow (kg) & 0.5302 & 0.2111 \\
Total external removal (kg) & 0.5302 & 0.7199 \\
\bottomrule\noalign{}
\end{tabular}
}
\endgroup
\end{table}

Both cases shared the same upper-wall film, an initially dry lower-wall film region and bulk flow held fixed. The 70.32\% lower-wall inventory reduction measures local storage, not separator efficiency. Added removal must include both exit paths: the drain removed 0.508790 kg, while film-edge outflow fell by 0.319124 kg. Net extra removal was 0.189666 kg, close to the 0.189694 kg decrease in total film mass. The difference, approximately 0.0000283 kg (0.015\% of extra removal), confirms close agreement between added removal and reduced storage.

In the longer Drain ON continuation, totals after the first Courant-limit crossing are excluded from performance interpretation; inner residuals were unavailable. Storage was already 37.35 kg/s over the preceding 15 ms. Upper-wall accumulation was therefore present before the numerical failure, locating the remaining storage problem above the drain.

\subsection{Wall resolution and mesh sensitivity}\label{a.6.-wall-resolution-and-mesh-sensitivity}

\begin{table}[!htbp]
\centering
\renewcommand{\thetable}{A7}
\caption{Main-wall mean \(y^+\) on the 60k mesh, supporting the wall-treatment assessment in Section 4.6. Each range spans two solver estimates of the area-weighted mean, not the variation across the wall.}\label{tab:A7}
\begingroup\small\setlength{\tabcolsep}{3pt}
{
\begin{tabular}{@{}
  >{\raggedright\arraybackslash}p{(\linewidth - 2\tabcolsep) * \real{0.5000}}
  >{\raggedleft\arraybackslash}p{(\linewidth - 2\tabcolsep) * \real{0.5000}}@{}}
\toprule\noalign{}
\begin{minipage}[b]{\linewidth}\raggedright
Comparison state
\end{minipage} & \begin{minipage}[b]{\linewidth}\raggedleft
Area-weighted mean \(y^+\): range of solver estimates
\end{minipage} \\
\midrule\noalign{}
Full-feed control, iteration 5586 & 29,331.95--34,334.64 \\
Combined treatment with contact-based absorber, iteration 45606 & 19,023.87--20,427.87 \\
Shortened startup, iteration 5080 & 18,757.91--20,140.20 \\
\bottomrule\noalign{}
\end{tabular}
}
\endgroup
\end{table}

Means use \(\sum A_jq_j/\sum A_j\), where \(A_j\) is face area and \(q_j\) is the local \(y^+\) value, over the same 3463 main-wall faces covering 53.436952 m². The mean distance from the wall to adjacent cell centres was 13.0089 mm; this is not the first-layer height. Layer count, first-layer height, orthogonal quality and skewness were not quantified in this assessment. Both \(y^+\) estimates are high, so the visible wall layers alone do not establish adequate near-wall resolution.

\begin{table}[!htbp]
\centering
\renewcommand{\thetable}{A8a}
\caption{Film storage and direct drainage at 195 ms, supporting Section 4.10.}\label{tab:A8a}
\begingroup\small\setlength{\tabcolsep}{3pt}
{
\begin{tabular}{@{}
  >{\raggedleft\arraybackslash}p{(\linewidth - 6\tabcolsep) * \real{0.2500}}
  >{\raggedleft\arraybackslash}p{(\linewidth - 6\tabcolsep) * \real{0.2500}}
  >{\raggedleft\arraybackslash}p{(\linewidth - 6\tabcolsep) * \real{0.2500}}
  >{\raggedleft\arraybackslash}p{(\linewidth - 6\tabcolsep) * \real{0.2500}}@{}}
\toprule\noalign{}
\begin{minipage}[b]{\linewidth}\raggedleft
Cells
\end{minipage} & \begin{minipage}[b]{\linewidth}\raggedleft
Total film (kg)
\end{minipage} & \begin{minipage}[b]{\linewidth}\raggedleft
Lower-wall film (g)
\end{minipage} & \begin{minipage}[b]{\linewidth}\raggedleft
Instantaneous direct drain (kg/s)
\end{minipage} \\
\midrule\noalign{}
60,964 & 7.490 & 34.07 & 22.716 \\
342,609 & 2.821 & 10.75 & 7.168 \\
679,970 & 2.950 & 10.17 & 6.778 \\
997,604 & 2.056 & 6.88 & 4.587 \\
\bottomrule\noalign{}
\end{tabular}
}
\endgroup
\end{table}

The drain law gives \(\dot m_D=M_{f,\mathrm{lower}}/\tau_f\): for example, \(0.03407/0.0015\simeq22.72\) kg/s on 60k. This is an instantaneous rate, not cumulative collection. Relative to 680k, the 997k film inventory and drain rate were 30.3\% and 32.3\% lower, using unrounded values.

Each film developed with its own mesh's bulk flow held fixed. Starting film times ranged from 0.700 to 0.900 ms, and the 342k case also included a further adjustment of the bulk flow. The 2.6M mesh supplied only an initial bulk-flow comparison and was excluded from the developed-film comparison. Table A8a therefore compares developing films under different starting conditions. These differences prevent the changes from being attributed to mesh resolution alone.

The bulk-flow comparison in Section 4.10 uses the same four 195 ms states. Static pressure difference is the area-averaged mixture pressure over both inlet faces minus that at the steam outlet. Outlet steam quality is \(100|\dot m_V|/(|\dot m_V|+|\dot m_L|)\), calculated from the outward phase mass-flow rates. These values describe the bulk flow held fixed during film development and exclude separately tracked droplets.

\begin{table}[!htbp]
\centering
\renewcommand{\thetable}{A8b}
\caption{Bulk-flow quantities at the 195 ms comparison points in Table A8a.}\label{tab:A8b}
\begingroup\small\setlength{\tabcolsep}{3pt}
{
\begin{tabular}{@{}
  >{\raggedright\arraybackslash}p{(\linewidth - 8\tabcolsep) * \real{0.2000}}
  >{\raggedleft\arraybackslash}p{(\linewidth - 8\tabcolsep) * \real{0.2000}}
  >{\raggedleft\arraybackslash}p{(\linewidth - 8\tabcolsep) * \real{0.2000}}
  >{\raggedleft\arraybackslash}p{(\linewidth - 8\tabcolsep) * \real{0.2000}}
  >{\raggedleft\arraybackslash}p{(\linewidth - 8\tabcolsep) * \real{0.2000}}@{}}
\toprule\noalign{}
\begin{minipage}[b]{\linewidth}\raggedright
Mesh
\end{minipage} & \begin{minipage}[b]{\linewidth}\raggedleft
Static pressure difference (kPa)
\end{minipage} & \begin{minipage}[b]{\linewidth}\raggedleft
Outlet vapour flow (kg/s)
\end{minipage} & \begin{minipage}[b]{\linewidth}\raggedleft
Outlet liquid flow (kg/s)
\end{minipage} & \begin{minipage}[b]{\linewidth}\raggedleft
Outlet steam quality (\%)
\end{minipage} \\
\midrule\noalign{}
60k & 11.0443 & 80.599794 & 0.548691 & 99.3238 \\
342k & 10.9847 & 81.172387 & 0.495744 & 99.3930 \\
680k & 11.0297 & 81.165782 & 0.496423 & 99.3921 \\
997k & 11.3209 & 81.112083 & 0.477213 & 99.4151 \\
\bottomrule\noalign{}
\end{tabular}
}
\endgroup
\end{table}

The values are evaluated at the comparison point, rather than averaged over time. The pressure range divided by its mean is 3.03\%. At these states, pressure and outlet steam quality varied much less than film storage and drainage.

For context, Pointon's typical predicted pressure-loss range starts around 10--15 kPa and extends to 40 kPa (p.~944). The 19 kPa pressure drop and 99.87 wt\% dryness in Table 1 are published design parameters, not measurements of the present separator (\href{https://publications.mygeoenergynow.org/grc/1028587.pdf}{Pointon et al., 2009, Table 1}).

\subsection{Evidence for simplifying the collection model}\label{a.7.-evidence-for-simplifying-the-collection-model}

\par\noindent\begin{minipage}{\linewidth}
The full-geometry trials support the choice in Section 4.3 to separate liquid transport from pool control. Six of 16 steady outlet cases reached 500 iterations, but completed cases still had large liquid mass-balance errors. Figure A1 shows the additional pressure-feedback test: liquid inventory kept rising as the pressure controller reached its lower bound.

\medskip
\centering
\includegraphics[width=0.98\linewidth,height=0.36\textheight,keepaspectratio]{\detokenize{figures/appendix-pool-control-history.png}}
\captionsetup{hypcap=false}\captionof{figure}[Pool inventory and pressure control]{Pool inventory and bounded outlet-pressure control.}
\label{fig:A1}
\end{minipage}\par\medskip

Liquid in the region \(y\leq0.10\) m increased from 187.79 to 284.83 kg during the 10,000-update trial and preceding 50-update check, while pressure reached its 1.115 MPa gauge lower bound at 98 of 100 control endpoints. The assumed 200 kg target was therefore not maintained with these control settings; this does not establish a limit on physical pool control. The result supports using a separate absorber to study transport towards the base without also solving pool and discharge control.

\subsection{Design evidence and its scope}\label{a.8.-design-evidence-and-its-scope}

The design tests in Section 4.11 explored local flow changes around the proposed features. The vertical blocker used 154,063 cells and the angled ramp 213,232 cells. Both basic-film cases reached 150 ms; their 50--150 ms continuations used 10,000 fixed 10 µs updates with bulk flow held fixed, stripping off and edge separation off. Different meshes, initial fields and preparation prevent an isolated geometry comparison.

Figure 21 presents the mixture-velocity plots associated with the 50 ms design cases. Viewing angles and vector sampling differ, so the comparison is qualitative. The vectors show bulk flow rather than film transport and cannot establish improved collection. They support Section 5.5 by showing where each feature changes local flow, giving longer film tests a specific effect to assess.

\end{document}
